\pdfoutput=1
\documentclass[11pt]{amsart}
\usepackage{amsmath,amssymb,amsthm}
\usepackage{mathtools}
\usepackage{mathrsfs}
\usepackage[margin=1.2in]{geometry}
\usepackage{booktabs}
\usepackage{array}
\usepackage{flafter}
\usepackage{xcolor}
\usepackage[colorlinks=true,linkcolor=blue,citecolor=blue,urlcolor=blue]{hyperref}
\hypersetup{pdfauthor={Bernd Johannes Wuebben},
  pdftitle={Point operations and matching-line bubbles in the Atiyah-Floer comparison},
  pdfsubject={Corrected point-operation compatibility in the admissible SO(3) Atiyah-Floer comparison},
  pdfkeywords={Atiyah-Floer conjecture, instanton Floer homology,
  Lagrangian Floer homology, operation algebra, mixed equation,
  matching-line bubbling, stable bundles, Hermitian Yang-Mills equation}}

\setlength{\emergencystretch}{2em}

\newtheorem{theorem}{Theorem}[section]
\newtheorem{proposition}[theorem]{Proposition}
\newtheorem{lemma}[theorem]{Lemma}
\newtheorem{corollary}[theorem]{Corollary}
\theoremstyle{definition}
\newtheorem{definition}[theorem]{Definition}
\newtheorem{convention}[theorem]{Convention}
\theoremstyle{remark}
\newtheorem{remark}[theorem]{Remark}

\theoremstyle{plain}
\newtheorem{letterresult}{Theorem}
\renewcommand{\theletterresult}{\Alph{letterresult}}
\newtheorem{corletter}[letterresult]{Corollary}
\renewcommand{\thecorletter}{\Alph{letterresult}}

\newcommand{\ZZ}{\mathbb{Z}}
\newcommand{\QQ}{\mathbb{Q}}
\newcommand{\RR}{\mathbb{R}}
\newcommand{\CC}{\mathbb{C}}
\newcommand{\PP}{\mathbb{P}}
\newcommand{\cA}{\mathcal{A}}
\newcommand{\cB}{\mathcal{B}}
\newcommand{\cD}{\mathcal{D}}
\newcommand{\cE}{\mathcal{E}}
\newcommand{\cF}{\mathcal{F}}
\newcommand{\cG}{\mathcal{G}}
\newcommand{\cL}{\mathcal{L}}
\newcommand{\cM}{\mathcal{M}}
\newcommand{\cN}{\mathcal{N}}
\newcommand{\cO}{\mathcal{O}}
\newcommand{\cP}{\mathcal{P}}
\newcommand{\cR}{\mathcal{R}}
\newcommand{\cS}{\mathcal{S}}
\newcommand{\frakM}{\mathfrak{M}}
\newcommand{\ad}{\operatorname{ad}}
\newcommand{\End}{\operatorname{End}}
\newcommand{\Hom}{\operatorname{Hom}}
\newcommand{\ev}{\operatorname{ev}}
\newcommand{\Ev}{\operatorname{Ev}}
\newcommand{\ind}{\operatorname{ind}}
\newcommand{\rank}{\operatorname{rank}}
\newcommand{\tr}{\operatorname{tr}}
\newcommand{\codim}{\operatorname{codim}}
\newcommand{\Glue}{\operatorname{Glue}}
\newcommand{\CS}{\operatorname{CS}}
\newcommand{\id}{\operatorname{id}}

\title[Point operations in the Atiyah--Floer comparison]{
Point operations and matching-line bubbles\\
in the Atiyah--Floer comparison}
\author{Bernd Johannes Wuebben}

\date{31 August 2026}
\subjclass[2020]{57R58 (primary); 53D40, 58J32, 14D20 (secondary)}
\keywords{Atiyah--Floer conjecture, instanton Floer homology,
Lagrangian Floer homology, operation algebra, mixed equation,
matching-line bubbling, stable bundles, Hermitian Yang--Mills equation}

\begin{document}

\begin{abstract}
Daemi, Fukaya, and Lipyanskiy proved the admissible $SO(3)$
Atiyah--Floer comparison and compatibility of the comparison map with the
operations associated with $H_1$ and $H_2$ of the splitting surface.  The
point class is the remaining generator; its comparison homotopy first
encounters energy concentration on the matching line, the codimension-one
seam between the gauge-theoretic and symplectic regions, in Fredholm index
four.

We compactify the one-dimensional point-cut mixed moduli space and classify
its new boundary.  Chern--Simons excision gives positive half-integral local
charge, which at index four forces a single charge-one matching bubble.
Relative stable reduction identifies the primitive limits, a Dirichlet
Hermitian--Yang--Mills problem reconstructs them, and a scale-uniform inverse
gives converse gluing.  Determinant-line excision reduces the oriented local
count to Mu\~noz's extension spaces.

The resulting boundary identity shows that uncorrected point-operation
compatibility fails in genera one and two.  After dividing the positive
Pontryagin point class by four, the symplectic point operation must be
shifted by $2$, $1$, or $0$ when the marked point lies on a component of
genus one, two, or at least three; in the integral $p_1$ normalization the
corresponding local counts are $8$, $4$, and $0$.  Together with the surface
and loop results of Daemi, Fukaya, and Lipyanskiy, this proves the full
corrected observable-algebra action in the admissible embedded-Lagrangian
setting.
\end{abstract}

\maketitle

\section{Introduction}\label{sec:intro}

\subsection{The point operation}
Let $(Y_\#,E_\#)$ be an admissible $SO(3)$ pair split along a closed
oriented surface $\Sigma$.  Atiyah's gauge--symplectic comparison
principle predicts an identification of the two resulting Floer theories
\cite{Atiyah88}.  In the embedded-Lagrangian setting of
Daemi--Fukaya--Lipyanskiy, henceforth DFL, the splitting produces an
instanton Floer complex $C_G$ with differential $d_G$, a Lagrangian Floer
complex $C_S$ with differential $d_S$, and a comparison chain map
\[
                     N:C_G\longrightarrow C_S.
\]
Their main theorem identifies the induced homologies
\cite[Theorem~4]{DaeFukLip21a}.  The comparison is defined by counting
solutions of a mixed equation: an anti-self-dual connection on one side
of a codimension-one seam and a pseudoholomorphic map on the other, with
their boundary values identified in the symplectic moduli space of flat
connections \cite{AtiyahBott83}.
The regularity, Fredholm, compactness, and removal theorems for this
equation are proved in \cite{DaeFukLip21b}.

The universal bundle over the flat moduli space and the splitting surface
gives operations indexed by $H_*(\Sigma)$.  DFL prove that $N$
intertwines the operations associated with $H_2(\Sigma)$ and
$H_1(\Sigma)$ \cite[Section~7.2]{DaeFukLip21a}.  They also explain why
the same argument stops at $H_0(\Sigma)$: a point insertion requires
four-dimensional mixed moduli spaces, and these can bubble on the
matching line.  A neighborhood of that boundary was not part of their
compactification.

The purpose of this paper is to isolate the obstruction and construct the
compactification needed to describe and count that neighborhood.  The
cylinder calculation shows that uncorrected compatibility is impossible in
low genus.  The general formula attributes the discrepancy to a
matching-line scalar depending only on the genus of the component carrying
the point.

The local charge, primitive stable reduction, Dirichlet reconstruction,
uniform inverse, no-neck theorem, primitive gluing collar, and stratified
point incidence are established below.  Determinant excision identifies
the local algebraic multiplicities and cylinder calibration fixes their
coherent signs.  The resulting oriented boundary formula proves the point
comparison and its full-operation consequence.

\begin{convention}\label{conv:global}
All moduli spaces use DFL's orientation, grading, and end conventions.
The symbol $p_1$ denotes the positive integral degeneracy cycle determined
by their product orientation.  Unless an integral statement is explicitly
made, coefficients are rational.  A \emph{matching atom} is an atom of the
mixed energy measure at a point of the matching line.  It is
\emph{primitive} when its relative clutching number is one, equivalently
when its normalized energy is $\frac12$.  Normalized energy means DFL's
topological energy, recalled in \eqref{eq:top-energy}.
\end{convention}

\subsection{Results}
Write
\[
                       \Sigma=\bigsqcup_j\Sigma_j,
       \qquad g_j=\operatorname{genus}(\Sigma_j).
\]
We retain all admissibility, irreducibility, and embeddedness hypotheses
of DFL.  In particular, the standard secondary perturbations are supported
away from the matching line.  For $q\in\Sigma_j$, let $b_q^G$ and $b_q^S$
be the gauge and symplectic operations represented by the positive
integral degeneracy cycle $p_1/[q]$.  Over $\QQ$ put
\[
                      m_q^G=\frac14b_q^G,
              \qquad m_q^S=\frac14b_q^S.
\]
This is DFL's positive quarter-Pontryagin convention.  The sign is fixed
by their product orientation on the linear-dependence cycle, rather than
by the alternative convention $\mu(q)=-\frac14p_1/[q]$ sometimes used in
instanton theory \cite{DonaldsonKronheimer90}.
For the cylinder splitting of $\Sigma_g\times S^1$, abbreviate the
corresponding point operations by $b_G$ and $b_S$.

\begin{letterresult}[low-genus cylinder obstruction;
  $=$ Proposition~\ref{prop:cylinder-identities}]
\label{thm:intro-cylinder}
For the cylinder splitting of $\Sigma_g\times S^1$, the positive integral
point operations satisfy
\[
 Nb_G-b_SN=
 \begin{cases}
 8N,&g=1,\\
 4N,&g=2.
 \end{cases}
\]
Consequently, the uncorrected point comparison cannot hold in general in
genus one or two.
\end{letterresult}

Theorem~A is proved independently of the matching-line compactification:
it follows from DFL's surface-operation comparison and Mu\~noz's cylinder
relations.  Proposition~\ref{prop:cylinder-identities} gives the proof.

\begin{letterresult}[point-class comparison; $=$ Theorem~\ref{thm:main}]
\label{thm:intro-main}
There are compatible secondary perturbations and a point-cut chain
homotopy $K_q:C_G\to C_S$ for which
\begin{equation}\label{eq:intro-main}
 N m_q^G-m_q^S N
 =d_SK_q+K_qd_G+\kappa_{g_j}N,
 \qquad
 \kappa_g=
 \begin{cases}
 2,&g=1,\\
 1,&g=2,\\
 0,&g\geq3.
 \end{cases}
\end{equation}
Equivalently, $N$ intertwines the gauge point operation with the corrected
symplectic operation
\[
              \widehat m_q^S=m_q^S+\kappa_{g_j}\id
\]
up to chain homotopy.
\end{letterresult}

Theorem~B is restated as Theorem~\ref{thm:main}; its integral
normalization is Theorem~\ref{thm:integral}.

The integral version of \eqref{eq:intro-main} has coefficients
$8$, $4$, and $0$; see Theorem~\ref{thm:integral}.  We state the quarter-Pontryagin
formula over $\QQ$ because an integral refinement requires a separate
divisibility choice for the universal bundle.

Theorem~B completes the list of generator comparisons.  Give
$H_k(\Sigma;\ZZ)$ degree $4-k$, and let $\mathbb A_{p_1}(\Sigma;R)$ denote
the free graded-commutative algebra on
$H_*(\Sigma;\ZZ)\otimes_{\ZZ}R$.

\begin{corletter}[full corrected operation algebra;
  $=$ Corollary~\ref{thm:full-operation-algebra}]
\label{thm:intro-operation-algebra}
For every commutative ring $R$, scalar extension of the DFL comparison is
an isomorphism of $\mathbb A_{p_1}(\Sigma;R)$-modules after the symplectic
point generator is corrected by $8$, $4$, or $0$ in genera one, two, or at
least three.  If $4$ is invertible in $R$, the quarter-Pontryagin correction
is respectively $2$, $1$, or $0$.  The intertwining identity for every
chosen word is induced by an explicit signed chain homotopy.
\end{corletter}

Corollary~C is the formal generator-by-generator consequence proved in
Corollary~\ref{thm:full-operation-algebra}.  Its only new input beyond DFL
is Theorem~B.

The scalar term comes from the following compactification theorem.  Let
$M_g$ denote the moduli space of flat connections on the
nontrivial $SO(3)$ bundle over a connected genus-$g$ surface.  Under the
Narasimhan--Seshadri correspondence it is the moduli space of stable
rank-two bundles with fixed odd determinant \cite{NarSes65}.  Mu\~noz's
extension space and compact primitive bundle space are defined as follows.
Fix the determinant line $\Lambda$, and let $\cE_g$ be the rank-$g$ vector
bundle on $J(\Sigma_g)$ with fiber
$H^1(\Sigma_g,L^2\otimes\Lambda^{-1})$ over $L$.  Then
\begin{equation}\label{eq:intro-munoz-spaces}
 N_g=\PP(\cE_g^\vee)\longrightarrow J(\Sigma_g),
 \qquad
 \frakM_g=\PP((\cE_g\oplus\cE_g)^\vee)
       \longrightarrow J(\Sigma_g),
\end{equation}
and there is a natural map $i_g:N_g\to M_g$
\cite{Munoz99b}.  A point of a fiber of $\frakM_g$ is a pair of extension
vectors $[v_0:v_1]$.  Independent vectors give a minimal rational curve;
dependent vectors give a degree-one quasimap with one basepoint.

For the next statement, a \emph{primitive matching bubble} means a rescaled
limit with relative clutching charge one.  We call it \emph{full rank} when
$v_0$ and $v_1$ are independent and \emph{rank drop} when they are dependent.
An algebraic limit is \emph{admissible along the seam} when every surface
restriction over the compactified matching line is stable and any quasimap
basepoint lies on the gauge-theoretic hemisphere.

\begin{letterresult}[primitive compactification and gluing]
\label{thm:intro-compactification}
Let a sequence in an index-four mixed moduli space develop a
matching-line bubble while meeting the point incidence.  Then its residual
mixed solution has index zero and the following hold.
\begin{enumerate}
\item There is one matching concentration point, its relative clutching
charge is one, its normalized energy is $\frac12$, and every neck and
smaller scale carries zero energy.
\item After passage to a subsequence, finite-scale holomorphic doubling and
relative stable reduction produce a primitive algebraic limit in
$\frakM_g$.  It is either a
full-rank minimal line or a rank-drop bundle with one positive Hecke
modification.  Both attach to the residual solution through
$i_g(N_g)\subset M_g$.
\item Every algebraic limit admissible along the seam is realized by an exact
local mixed solution.  The primitive stratum has a gluing collar, and the
gluing map is a local diffeomorphism onto a neighborhood of the bubbled face.
\item After the integral $p_1$ point cut, the oriented local contribution
per residual solution is
\[
                     c_g=
       \begin{cases}
       8,&g=1,\\
       4,&g=2,\\
       0,&g\geq3.
       \end{cases}
\]
\end{enumerate}
\end{letterresult}

The atomwise charge statement in Theorem~D and
the uniqueness of the physical matching center follow from
Theorem~\ref{thm:local-charge} and
Corollary~\ref{cor:primitive-atom}.  Proposition~\ref{prop:primitive-no-neck}
excludes analytic neck energy, and Theorems~\ref{thm:reconstruction} and
\ref{thm:primitive-gluing} reconstruct the primitive bubble and give its
collar.  Proposition~\ref{prop:point-manifolds} supplies the transverse
compact point cut.  Proposition~\ref{prop:determinant-excision} and
Section~\ref{sec:coefficients} prove part~(4).

The two types in part~(2) are both necessary in the compactification.
Convergence away from a
physical concentration point fixes the rescaled evaluation at infinity;
it does not force a rescaled map to be constant at finite points.  The
generic genus-two bubble is therefore full rank, while rank drop is a
compactification stratum.  In genus one every pair of extension vectors is
dependent and only the rank-drop, purely gauge-theoretic branch remains.

\subsection{Why the correction is forced}
The low-genus constants can be seen before doing any matching-line
analysis.  Take the cylinder splitting of $\Sigma_g\times S^1$.  On the
gauge side the point operation is multiplication by Mu\~noz's degree-four
Floer generator $b_G$; on the symplectic side it is quantum multiplication
by the geometric generator $b_S$.  For $g=2$ the ring relations are
\[
                 a_G^2+b_G-8=0,
             \qquad a_S^2+b_S-4=0,
\]
where $a$ is the degree-two surface generator
\cite{Munoz99a,Munoz99b}.  DFL already prove $Na_G=a_SN$, and hence
\[
                      Nb_G-b_SN=4N.
\]
For $g=1$, the odd-determinant character variety is a point, whereas the
instanton Floer relation is $b_G=8\id$; therefore
\[
                      Nb_G-b_SN=8N.
\]
These identities use only the established surface compatibility and do
not assume point compatibility.  They prove that an uncorrected theorem
is impossible and later calibrate the local determinant orientation.

\subsection{Method}
The point-cut homotopy is zero-dimensional in mixed index three.  Its
boundary relation is one-dimensional in mixed index four.  DFL's
energy-index formula and half-integrality imply that any defective
index-four sequence loses exactly $\frac12$ unit of topological energy and
has an index-zero residual solution.  This global statement does not by
itself prevent several nonquantized matching atoms.

The proof has four stages.  First, on a small circle about a matching point,
the gauge and symplectic halves are spliced to a connection on
$S^1\times\Sigma_g$.  Comparing the original smooth filling with the
removable limiting filling identifies their mixed Chern--Simons branch
difference with $n/2$, where $n$ is the Pontryagin number of a relative
clutching bundle on $S^2\times\Sigma_g$.  The total defect $1/2$ then forces
$n=1$ at one point.  A good-circle Chern--Simons argument excludes energy
on the intervening neck.

The second step avoids the usual adiabatic projection.  After
base dilation the vertical-curvature term in the ASD equation is
multiplied by the square of the scale, rather than its inverse.  The usual
adiabatic projection to $M_g$ therefore does not apply; compare
\cite{DosSal94,DosSal07}.  Instead, every smooth finite-scale mixed
solution determines a holomorphic bundle on the full disc times $\Sigma$.
Relative stable reduction associates to these doubled bundles a classifying
map and a finite modification cycle.  Relative $p_1$ additivity gives
\[
                            n=d+\ell.
\]
The primitive identity $n=1$ leaves exactly a minimal line
$(d,\ell)=(1,0)$ or one rank-drop modification $(d,\ell)=(0,1)$.

The third stage is reconstruction and gluing.  Decomposing $\PP^1$ into
two hemispheres and prescribing the Narasimhan--Seshadri seam metrics puts
the gauge hemisphere in the setting of Donaldson's Dirichlet theorem for
Hermitian Yang--Mills metrics \cite{Donaldson92}.  The Dirichlet solution
reconstructs every seam-admissible algebraic point.  The reduced Fredholm
problem is the Dolbeault deformation complex, and an anisotropic
difference-quotient estimate keeps its right inverse bounded as the gauge
hemisphere collapses.  Pregluing and contraction then give the primitive
collar.

The fourth stage treats incidence and orientation.  Stratified
transversality extends the point cut over the collar.  Determinant-line
excision separates the residual orientation from a universal local scalar.
The image of $i_g$ has real codimension $2(g-2)$, the
genus-two map $i_2$ has degree four, and the genus-one excess line has degree
eight.  The Dolbeault determinant identifies these numbers with the
matching-face counts, and the cylinder identities fix both coherent signs
as positive.

\subsection{Relation to prior work}
The mixed equation and the admissible comparison theorem are due to DFL
\cite{DaeFukLip21a,DaeFukLip21b}.  The earlier adiabatic comparison of
instantons and holomorphic curves is due to Dostoglou--Salamon
\cite{DosSal94,DosSal07}.  Wehrheim established energy quantization and
removal for ASD instantons with Lagrangian boundary conditions
\cite{Wehrheim05a} and the energy identity on $\CC\times\Sigma$
\cite{Wehrheim05b}.  Owens proved the correspondence between finite-energy
instantons on a cylinder and stable bundles on $\PP^1\times\Sigma$
\cite{Owens01}; the stable-family-plus-modification description used here
is its relative local form.

Mu\~noz computed both the instanton Floer ring of $\Sigma\times S^1$ and
the quantum cohomology of $M_g$ \cite{Munoz99a,Munoz99b}.  His
projective-bundle description of primitive product bundles supplies the
finite-dimensional compactification and the two low-genus numbers.  The
appearance of a correction term is analogous in spirit to seam-bubble
corrections in pseudoholomorphic quilts \cite{BottmanWehrheim18}, but the
local object here couples an ASD half to a holomorphic half and is not
identified with a figure-eight bubble.

\section{The point cut in the mixed moduli space}
\label{sec:point-cut}

\subsection{The DFL setting}\label{sec:dfl-setting}
Let
\[
                      Y_\#=Y\cup_\Sigma(-Y')
\]
be an admissible splitting in the sense of
\cite[Sections~2--3]{DaeFukLip21a}.  The restriction $F$ of the
$SO(3)$ bundle to every component of $\Sigma$ is nontrivial.  We assume,
as in the embedded form of the DFL theorem, that the perturbed
three-manifold Lagrangians are embedded and intersect transversely and that
all flat connections in the construction are irreducible.  Put
\[
 M=\cM(\Sigma,F)=\prod_jM_j,
 \qquad M_j=\cM(\Sigma_j,F_j).
\]
The standard complex structure $J_*$ identifies $M_j$ with the
fixed-odd-determinant stable-bundle moduli space.  Its real dimension is
$6g_j-6$ for $g_j\geq2$; $M_j$ is a point for $g_j=1$.

For generators $\alpha\in C_G$ and $\beta\in C_S$, let
\[
                         \cM_\eta(\alpha,\beta)_r
\]
denote the moduli space of perturbed mixed solutions whose linearized
mixed operator has Fredholm index $r$.  The secondary perturbation $\eta$
is a finite sum of gauge-theoretic holonomy perturbations and
domain-dependent almost-complex variations.  It is identically zero on a
fixed neighborhood of the matching line.  DFL construct such
perturbations regular through index three
\cite[Proposition~3.23]{DaeFukLip21a}.

For a mixed pair $(A,u)$, their normalized topological energy is
\begin{equation}\label{eq:top-energy}
 \cE(A,u)=\frac1{8\pi^2}\int_X\tr(F_A\wedge F_A)
          +\frac1{4\pi^2}\int_Su^*\Omega,
\end{equation}
with the perturbation terms incorporated on the cylindrical pieces.  The
index formula has the form
\begin{equation}\label{eq:index-energy}
 \ind\cD_{(A,u)}
 =8\cE(A,u)+\varepsilon(\beta)-\varepsilon(\alpha).
\end{equation}
For two mixed pairs with the same ends,
\begin{equation}\label{eq:half-integrality}
 2\bigl(\cE(A,u)-\cE(A',u')\bigr)\in\ZZ;
\end{equation}
see \cite[Lemma~3.15]{DaeFukLip21a}.  These two formulas are the
numerical source of the index-four threshold.

\subsection{The universal degeneracy cycle}\label{sec:degeneracy}
Let $\widetilde{\mathbb V}$ be DFL's interpolating universal adjoint
bundle over the mixed configuration space times the glued parameter space
$\widetilde X$.  Its complexification has rank three.  For two generic
sections $s_1,s_2$, let
\[
 \mathscr D(s_1,s_2)
 =\{x:\rank_\CC(s_1(x),s_2(x))\leq1\}.
\]
At the top rank-one stratum this has complex codimension two and hence
real codimension four.  With DFL's product orientation it represents the
positive class $p_1(\widetilde{\mathbb V})$.  The rank-zero stratum has
larger codimension and is absent from every moduli space considered below.

The interpolation from the gauge insertion to the symplectic insertion is
an embedding
\[
                    \iota:\RR\times\Sigma\longrightarrow\widetilde X.
\]
The parameter $t\in\RR$ tends at the two ends to the representatives used
to define $b_q^G$ and $b_q^S$.  If $P\subset\Sigma$ is an oriented smooth
$k$-cycle, define
\[
 \cL_d^P(\alpha,\beta)
 \subset \cM_\eta(\alpha,\beta)_r\times\RR_t\times P
\]
by the condition
\[
 s_1([A,u],\iota(t,x))\wedge
 s_2([A,u],\iota(t,x))=0.
\]

\begin{lemma}[dimension of the comparison cut]\label{lem:cut-dimension}
If the mixed moduli space and the degeneracy cycle are transverse, then
$\cL_d^P(\alpha,\beta)$ has expected dimension $d$ precisely when
\begin{equation}\label{eq:cut-dimension}
                              r=d+3-k.
\end{equation}
\end{lemma}

\begin{proof}
The mixed moduli, the interpolation parameter, and the cycle $P$ have
total dimension $r+1+k$.  The top degeneracy stratum has real codimension
four.  Thus the cut has dimension $r+k-3$, which equals $d$ exactly when
\eqref{eq:cut-dimension} holds.
\end{proof}

For a point $q\in\Sigma$, the zero-dimensional comparison homotopy uses
index three and the one-dimensional relation uses index four:
\begin{equation}\label{eq:point-cuts}
 \begin{aligned}
 \cL_0^q(\alpha,\beta)
 &\subset\cM_\eta(\alpha,\beta)_3\times\RR_t,\\
 \cL_1^q(\alpha,\beta)
 &\subset\cM_\eta(\alpha,\beta)_4\times\RR_t.
 \end{aligned}
\end{equation}
This is the first insertion for which DFL's construction reaches the
matching-line bubbling threshold.

\subsection{The point-cut homotopy}\label{sec:candidate-homotopy}
By Proposition~\ref{prop:index-exhaustion} below, mixed moduli spaces of
index at most three do not bubble.  Once transversality is imposed, the
oriented zero-manifolds $\cL_0^q(\alpha,\beta)$ therefore define
\begin{equation}\label{eq:K-integral}
 K_q^{p_1}(\alpha)
 =\sum_\beta\#\cL_0^q(\alpha,\beta)\,\beta.
\end{equation}
Set $K_q=\frac14K_q^{p_1}$ over $\QQ$.

The ordinary compactification of $\cL_1^q$ has four familiar boundary
types: gauge and symplectic Floer breaking and escape of $t$ to its gauge
and symplectic ends.  With DFL's conventions their signed sum is
\begin{equation}\label{eq:ordinary-boundary}
 d_SK_q^{p_1}+K_q^{p_1}d_G
 -Nb_q^G+b_q^SN.
\end{equation}
The only possible additional codimension-one face consists of
matching-line bubbles.  We denote its signed contribution by $B_q$ until
Section~\ref{sec:coefficients}.  The desired integral boundary formula is
therefore
\begin{equation}\label{eq:boundary-with-B}
 Nb_q^G-b_q^SN
 =d_SK_q^{p_1}+K_q^{p_1}d_G+B_q.
\end{equation}

The remainder of the paper constructs this face and establishes
$B_q=c_{g(q)}N$ with $c_1=8$, $c_2=4$, and $c_g=0$ for $g\geq3$.

\section{Index-four concentration and local charge}
\label{sec:charge}

\subsection{Exhaustion of the index budget}\label{sec:index-exhaustion}
Consider a sequence in $\cM_\eta(\alpha,\beta)_r$.  After DFL chain
compactness, write $(A_0,u_0)$ for the residual mixed solution,
$A_i^G$ for the nonconstant gauge Floer levels, and $v_j^S$ for the
nonconstant symplectic Floer levels.  Define the global defect
\begin{equation}\label{eq:global-defect}
 \Delta=\cE(A_\nu,u_\nu)-\cE(A_0,u_0)
        -\sum_i\cE(A_i^G)-\sum_j\cE(v_j^S).
\end{equation}
It includes all energy not carried by the displayed residual levels:
ordinary bubbles, matching atoms, and any neck loss.

\begin{proposition}[index-four exhaustion]\label{prop:index-exhaustion}
Assume that the residual moduli spaces are regular and have no
negative-index solutions.
\begin{enumerate}
\item A sequence of mixed solutions of index at most three cannot bubble.
\item If an index-four sequence has $\Delta>0$, then
\begin{equation}\label{eq:exhaustion}
 \Delta=\frac12,\qquad
 \ind\cD_{(A_0,u_0)}=0,
\end{equation}
and there are no nonconstant gauge or symplectic Floer levels.
\item The defect in part~(2) carries all four missing Fredholm-index
units.
\end{enumerate}
\end{proposition}

\begin{proof}
DFL compactness makes $\Delta$ nonnegative.  If it is positive, the
strict form of the topological-energy inequality and
\eqref{eq:half-integrality} give $\Delta\geq\frac12$.  Additivity of the
mixed index under gluing gives
\begin{equation}\label{eq:index-ledger}
 \ind\cD_{(A_0,u_0)}+
 \sum_i\ind\cD_{A_i^G}+
 \sum_j\ind\cD_{v_j^S}+8\Delta\leq r.
\end{equation}
Every term except the last is nonnegative, and each nonconstant Floer
level has positive index.  If $r\leq3$, the inequalities
$8\Delta\geq4$ and \eqref{eq:index-ledger} are incompatible.  If $r=4$,
equality is forced throughout: $8\Delta=4$ and every residual index other
than the index-zero mixed component vanishes.  A nonconstant Floer level
cannot have index zero, proving the claim.
\end{proof}

An interior four-dimensional instanton has integral charge at least one,
hence energy at least one and index at least eight.  It cannot occur in
Proposition~\ref{prop:index-exhaustion}(2), either as the primary bubble or
as a smaller-scale bubble inside a matching configuration.  A
nonconstant sphere in $M_g$ has normalized energy at least $\frac12$ and
Chern index four; a nonconstant disc on either three-manifold Lagrangian
has Maslov index at least four.  Thus the half-unit budget permits at most
one ordinary symplectic bubble and no simultaneous breaking.

\subsection{Ordinary bubbles and the point incidence}
\label{sec:ordinary-bubbles}
The preceding numerical statement concerns the uncut moduli space.  The
point incidence removes the remaining ordinary sphere and disc strata from
the finite-$t$ boundary.

\begin{proposition}\label{prop:ordinary-bubbles}
For a generic stratified representative of the universal $p_1$ cycle, no
finite-$t$ boundary point of $\overline{\cL}_1^q(\alpha,\beta)$ contains an
interior sphere bubble or an outer-boundary disc bubble.
\end{proposition}

\begin{proof}
Put $2n=\dim_\RR M_g$.  The moduli space of simple minimal spheres with
one nodal marked point has virtual dimension
\[
                    (2n-6)+2c_1(A_g)+2=2n,
\]
because $c_1(A_g)=2$.  Constraining the nodal evaluation to the residual
map removes $2n$ dimensions.  The remaining two variables are the position
of the node in the symplectic domain.  Thus the stable-sphere stratum has
codimension two in the uncut index-four mixed moduli space.

If the point incidence remains on the residual component, it is evaluated
on the one-dimensional interpolation family of a fixed index-zero
residual solution.  A generic codimension-four cycle misses that family.
If the incidence enters the sphere scale, requiring the bubble point to
equal the point on the fixed interpolation path leaves one real domain
variable; a marked point on the sphere adds two, and the $p_1$ incidence
removes four.  The resulting virtual dimension is $1+2-4=-1$.

For an outer-boundary Maslov-four disc, the interpolation path is separated
from the boundary, so the incidence cannot enter the disc scale.  On the
residual component it again sees only a one-dimensional family and is
missed generically.  Stable spheres or discs occurring after
$t\to+\infty$ are part of the compactified symplectic point operation and
are already included in $b_q^SN$; they are not finite-$t$ correction
terms.
\end{proof}

The matching line is different.  The interpolation path crosses the seam,
so the ratio between its normal coordinate and the concentration scale is
a parameter in the rescaled problem.  Together with the center, scale,
and marked-domain directions, it provides the four real directions cut by
the incidence.  A matching bubble can therefore be a genuine boundary
point of the one-dimensional cut.

\subsection{The local Chern--Simons sector}\label{sec:local-CS}
DFL's half-integrality controls the total defect, not each atom.  The next
theorem is the local quantization needed to rule out several matching
centers whose a priori energy measures sum to $\frac12$.

Let $z$ be a matching point at which a sequence $(A_\nu,u_\nu)$
concentrates.  Choose a small half-disc on either side of the seam, with
no other concentration point in its closure.  On the negative semicircle
use the restriction of $A_\nu$.  On the positive semicircle choose the
flat representatives of $u_\nu$ supplied by the matching condition.  At
the two seam endpoints insert the standard flat cylinder.  These pieces
give a spliced connection
\[
                       B_{\nu,r}\quad\text{on}\quad
                       Y_r\cong S^1\times\Sigma_g.
\]
The removable limit gives a second connection $B_{0,r}$ on the same
three-manifold.

DFL's connection-valued filling construction replaces this boundary data
by a smooth connection on $D^2\times\Sigma_g$.  Its Chern--Weil integral
contains both the gauge energy and the symplectic area
\cite[Lemma~4.67]{DaeFukLip21b}; the equality between the curvature term of
the flat-representative family and twice its symplectic area is
\cite[Lemma~5.10]{DaeFukLip21a}.

\begin{theorem}[local charge]\label{thm:local-charge}
Let $\delta_z>0$ be the atom of the normalized mixed energy measure at
$z$.  The original bundle and the removable limiting bundle determine a
relative clutching integer $n_z$ such that
\begin{equation}\label{eq:local-charge}
                        \delta_z=\frac{n_z}{2},
                 \qquad n_z\in\ZZ_{>0}.
\end{equation}
The sign is normalized so that the DFL positive generator has $n_z=1$ and
$p_1=-2$.
\end{theorem}

\begin{proof}
Take regular radii $r\downarrow0$ whose closed discs contain no other
concentration point.  For fixed $r$, compactness supplies annular gauges in
which the boundary connections converge smoothly.  Apply DFL's filling
construction to the original and removable limiting pairs, and denote the
resulting bundles with connection over $D^2\times\Sigma_g$ by
$(Q_{\nu,r},\mathbb A_{\nu,r})$ and $(Q_{0,r},\mathbb A_{0,r})$.  Use the
limiting framing to identify $Q_{0,r}|_{Y_r}$ with a fixed bundle
$Q_r^\partial\to Y_r$.  The annular convergence gauge supplies a projective
bundle isomorphism
\[
 \gamma_{\nu,r}:Q_{\nu,r}|_{Y_r}\longrightarrow Q_r^\partial
\]
for which $\gamma_{\nu,r*}B_{\nu,r}\to B_{0,r}$ smoothly.  No extension of
$\gamma_{\nu,r}$ over either filling is chosen or required; a nonextendible
gauge class will be the clutching class below.  With
DFL's normalization, the filling identities just cited are
\[
 \cE_\nu(r)=\frac{1}{8\pi^2}
   \int_{D^2\times\Sigma_g}\tr(F_{\mathbb A_{\nu,r}}\wedge
                                      F_{\mathbb A_{\nu,r}}),
 \qquad
 \cE_0(r)=\frac{1}{8\pi^2}
   \int_{D^2\times\Sigma_g}\tr(F_{\mathbb A_{0,r}}\wedge
                                      F_{\mathbb A_{0,r}}).
\]

Interpolate linearly, in a fixed affine chart of connections on
$Q_r^\partial$, between $\gamma_{\nu,r*}B_{\nu,r}$ and $B_{0,r}$ on
$[0,1]\times Y_r$.  Smooth convergence makes its normalized Chern--Weil
integral $\epsilon_{\nu,r}$ tend to zero.  Glue the original filling through
$\gamma_{\nu,r}$, this interpolation, and the reverse of the limiting
filling.  The result is an $SO(3)$ bundle $P_{\nu,r}$ over
$S^2\times\Sigma_g$.  Thus every component of the annular gauge, including
a possible large-gauge component, is recorded in $P_{\nu,r}$.  Replacing
$\gamma_{\nu,r}$ by a convergence gauge in the same local slice changes the
clutching map by a homotopy and leaves $p_1(P_{\nu,r})$ unchanged.
Chern--Weil excision,
with the sign convention of \cite[Proposition~5.14]{DaeFukLip21a}, gives
\begin{equation}\label{eq:relative-p1}
 \cE_\nu(r)-\cE_0(r)+\epsilon_{\nu,r}
       =-\frac14p_1(P_{\nu,r}).
\end{equation}

The manifold $S^2\times\Sigma_g$ is spin.  The Wu formula and
$p_1(P)\equiv w_2(P)^2\pmod2$ therefore show that $p_1(P_{\nu,r})$ is
even.  Put
\[
                    n_{\nu,r}:=-\frac12p_1(P_{\nu,r})\in\ZZ.
\]
For fixed $r$, weak convergence of the energy measures and
$\epsilon_{\nu,r}\to0$ show that the half-integers $n_{\nu,r}/2$ converge.
They are consequently eventually constant; write their limit as $n(r)/2$.
As regular radii tend to zero,
\[
 \frac{n(r)}2
 =\lim_{\nu\to\infty}\bigl(\cE_\nu(r)-\cE_0(r)\bigr)
 \longrightarrow\delta_z,
\]
because the removable limit is smooth.  The convergent half-integers
$n(r)/2$ are again eventually constant.  Calling the eventual integer
$n_z$ proves \eqref{eq:local-charge}.  Positivity of the atom gives
$n_z>0$.  DFL's disc generator in Section~2.2 has normalized curvature
integral $1/2$, so it has $p_1=-2$ and fixes the stated sign.

All Chern--Simons lifts in this argument are the real-valued lifts selected
by the displayed fillings.  For two regular radii, restricting the original
and limiting fillings to the intervening annulus gives the same clutching
bundle on their common boundary.  Hence changing the radius cannot transfer
an integer to an independently chosen boundary branch; the only change is
the Chern--Weil integral of that actual annular restriction, already present
in \eqref{eq:relative-p1}.
\end{proof}

The construction does not pass through $M_g$.  It therefore retains the
clutching class when $g=1$, even though $M_1$ is a point.  It is also
additive over disconnected components of $\Sigma$, since the relative
Pontryagin number splits componentwise.

\begin{corollary}[unique primitive matching center]\label{cor:primitive-atom}
Suppose an index-four sequence contributing to the point cut develops a
matching atom.  Then there is exactly one matching concentration point
$z$, and
\begin{equation}\label{eq:primitive-charge}
                 n_z=1,\qquad \delta_z=\Delta=\frac12.
\end{equation}
There is no second matching point and no ordinary bubble or nonconstant
Floer level away from $z$.
\end{corollary}

\begin{proof}
Proposition~\ref{prop:index-exhaustion} gives total defect $\frac12$.
Theorem~\ref{thm:local-charge} assigns every matching atom a positive
multiple of $\frac12$.  Thus the first atom already carries the whole
defect and has charge one.  Every ordinary sphere or disc also costs at
least $\frac12$, and every interior instanton costs at least one.  Energy
exhaustion excludes each such component away from the matching center.
\end{proof}

Combining Proposition~\ref{prop:ordinary-bubbles} and
Corollary~\ref{cor:primitive-atom}, the only new finite-$t$ concentration
locus is one charge-one matching center attached to an index-zero residual
mixed solution.  Quantization of the atom does not by itself exclude nested
scales or neck energy at that same center.  The next two sections prove the
stable-reduction and gluing results needed to do so.

\section{Holomorphic compactification of a primitive atom}
\label{sec:stable-reduction}

\subsection{The local equations and base dilation}\label{sec:local-equations}
Choose coordinates $(s,\theta)$ centered at the matching point, with
$D_-=\{s\leq0\}$ the gauge half-disc and $D_+=\{s\geq0\}$ the symplectic
half-disc.  In a gauge adapted to the $s$ direction, write the connection
on $D_-\times\Sigma_g$ as
\[
                         A=\beta(s,\theta)+\psi(s,\theta)d\theta.
\]
The ASD equation is
\begin{align}
 *_\Sigma F_\beta+\partial_s\psi&=0,
                                      \label{eq:local-asd-1}\\
 -\partial_\theta\beta+*_\Sigma\partial_s\beta
                  +d_\beta\psi&=0.  \label{eq:local-asd-2}
\end{align}
On $D_+$ the map $u$ satisfies
\begin{equation}\label{eq:local-CR}
                   \partial_\theta u-J_*\partial_su=0.
\end{equation}
Along $s=0$, the connection $\beta(0,\theta)$ is flat and represents
$u(0,\theta)$.

For a concentration scale $\rho\to0$, set
\[
 \beta_\rho(S,\Theta)=\beta(\rho S,\rho\Theta),
 \qquad
 \psi_\rho(S,\Theta)=\rho\psi(\rho S,\rho\Theta).
\]
Equations \eqref{eq:local-asd-1}--\eqref{eq:local-asd-2} become
\begin{align}
 \rho^2 *_\Sigma F_{\beta_\rho}+\partial_S\psi_\rho&=0,
                                      \label{eq:scaled-asd-1}\\
 -\partial_\Theta\beta_\rho+*_\Sigma\partial_S\beta_\rho
                  +d_{\beta_\rho}\psi_\rho&=0.  \label{eq:scaled-asd-2}
\end{align}
The coefficient of vertical curvature is $\rho^2$, not $\rho^{-2}$.
This is the opposite of the Dostoglou--Salamon adiabatic regime and does
not force the interior surface slices to be flat.  The correct intermediate
object is a holomorphic bundle on the product, rather than a direct map of
the negative half-plane to $M_g$.

\subsection{Finite-scale doubling}\label{sec:doubling}
Give the base the complex coordinate $w=s+i\theta$.

\begin{lemma}[finite-scale holomorphic doubling]\label{lem:doubling}
Every smooth mixed solution on a disc neighborhood of the matching line
determines a holomorphic rank-two bundle on the full disc times $\Sigma_g$.
It is canonical up to fixed-determinant complex gauge and tensoring by a
line bundle pulled back from the disc.
\end{lemma}

\begin{proof}
On the negative half, the product metric is K\"ahler.  The ASD equation
implies $F_A^{0,2}=0$, so $A^{0,1}$ is an integrable Dolbeault operator.

Every seam fiber is stable.  Stability is open, so the negative
holomorphic family has stable fibers on a collar of the seam.  The fine
moduli property \cite[Section~4]{Owens01} gives a holomorphic classifying
map $f_-$ on that collar and identifies its bundle, up to a line bundle
from the base, with the pullback of the universal bundle.  The positive
half similarly has the holomorphic classifying map $f_+=u$.

The matching condition identifies the continuous boundary values of
$f_-$ and $f_+$ on the seam.  In a complex chart of $M_g$, their union is
continuous on a full collar and holomorphic on its two open halves.
Morera's theorem shows that it is holomorphic across the seam.  Pull back
the universal bundle by this glued map and use the resulting collar
isomorphisms as transition maps between the negative and positive
families.  This gives the required holomorphic bundle.  The fine-moduli
ambiguity is a line bundle from the base; changing its trivialization and
the collar isomorphisms has exactly the ambiguity stated in the lemma.
\end{proof}

The point of Lemma~\ref{lem:doubling} is that it is applied before taking
the rescaled limit.  No unproved trace convergence for a limiting
connection-valued half-plane equation is required.

\subsection{Stable reduction over a disc}\label{sec:disc-reduction}
The following one-parameter statement is the local stable-reduction input.
It is the bundle-theoretic mechanism appearing in Owens's cylindrical
correspondence \cite{Owens01}.

\begin{lemma}[disc stable reduction]\label{lem:disc-reduction}
Let $E\to D\times\Sigma_g$ be a holomorphic rank-two bundle whose surface
fibers have fixed determinant of degree one, and suppose at least one fiber
is stable.
Then:
\begin{enumerate}
\item the unstable fibers form a finite analytic subset of every compact
subdisc;
\item the classifying map from the stable locus to $M_g$ extends
holomorphically across the unstable fibers without base change;
\item the original and pullback stable families differ by a finite Langton
sequence supported on those fibers; all terms are locally free, retain the
fixed surface-fiber determinant, and agree with $E$ off those fibers;
\item if one step is the kernel of a quotient line $G_j$ of degree
$q_j\leq0$, its positive charge is
\[
 \ell_j=\frac12\bigl(p_1(\ad E^{j+1})-p_1(\ad E^j)\bigr)
       =1-2q_j>0.
\]
Thus $\ell=\sum_j\ell_j$ is a nonnegative integer, and it vanishes exactly
when no modification occurs.
\end{enumerate}
\end{lemma}

\begin{proof}
The family $E$ is flat over $D$.  Toma's relative
Harder--Narasimhan theorem makes semistability Zariski open in such an
analytic family \cite[Theorem~2.4]{Toma20}.  Here semistability is stability,
so the existence of one stable fiber makes the unstable locus a proper
analytic subset of $D$, hence finite on compact subdiscs.

Fix an unstable point and a smaller disc $\Delta$ containing no other one.
Toma's analytic Langton theorem produces, without base change, a coherent
subsheaf $E'\subset E$ that equals $E$ over $\Delta^*$, is flat over
$\Delta$, and has stable central fiber
\cite[Theorem~3.1]{Toma20}.  It is locally free.  Indeed, all its fibers are
torsion-free sheaves on the smooth curve and hence are vector bundles.  At a
stalk with base parameter $t$, lift a basis of $E'/tE'$ to a map
$R^2\to E'$.  Nakayama makes it surjective; flatness gives
$\operatorname{Tor}_1^R(E',R/(t))=0$, so Nakayama applied to the kernel
makes the map an isomorphism.

In rank two Toma's construction is the finite Langton sequence
\[
0\longrightarrow E^{j+1}\longrightarrow E^j
 \longrightarrow i_*G_j\longrightarrow0,
\]
where, on the central curve, $B_j\subset E^j_0$ is the saturated
destabilizing line and $G_j=E^j_0/B_j$.  The transformed central fiber fits
into
\[
0\longrightarrow G_j\longrightarrow E^{j+1}_0
 \longrightarrow B_j\longrightarrow0.
\]
Thus every step retains the determinant $B_j\otimes G_j$ and equals its
predecessor off the central fiber.  The final stable family supplies the
holomorphic extension of the classifying map by the fine-moduli property;
it differs from the pullback universal family only by a line bundle from
$\Delta$.  In particular, the comparison preserves the outer attaching
identification.  This is also the local data in
\cite[Proposition~5.2]{Owens01}.

Put $q_j=\deg G_j$.  Since $\deg E^j_0=1$ and $B_j$ destabilizes,
$q_j\leq0$.  For the displayed elementary transform, the fiber divisor has
square zero and
\[
 c_1(E^{j+1})=c_1(E^j)-[\Sigma_g],\qquad
 \operatorname{ch}_2(E^{j+1})=\operatorname{ch}_2(E^j)-q_j.
\]
It follows that
\[
p_1(\ad E^{j+1})-p_1(\ad E^j)=2-4q_j.
\]
This proves the charge formula and positivity.  Repeating at the finitely
many unstable fibers proves the lemma.
\end{proof}

\subsection{The primitive charge splitting}\label{sec:charge-splitting}
Let $E\to D\times\Sigma_g$ be one doubled family supplied by
Lemma~\ref{lem:doubling}.  The convergence to the removable residual
solution fixes an outer projective framing.  Apply
Lemma~\ref{lem:disc-reduction}, denote its final stable family by $F$, its
classifying map by $f:D\to M_g$, and its total modification charge by
$\ell$.

Choose the outer circle so that $f(\partial D)$ lies in a contractible
fine-moduli chart.  Cap it in that chart and put
\begin{equation}\label{eq:capped-degree}
 d=\left\langle
 -\frac12p_1(\ad\mathcal U)/[\Sigma_g],
 [\widehat f(S^2)]\right\rangle .
\end{equation}
This is the positive generator normalization: the pullback of
$\ad\mathcal U$ by a minimal line has Pontryagin number $-2$
\cite[Proposition~3]{Munoz99b}.

\begin{proposition}[fixed-family framed splitting]
\label{prop:framed-splitting}
If $n(E)$ is the relative clutching charge in the outer framing, then
\begin{equation}\label{eq:finite-charge-splitting}
                              n(E)=d+\ell.
\end{equation}
The identity is independent of the cap, the universal-bundle
representative, the base-line ambiguity in the fine-moduli problem, and
the factorization into elementary transforms.  In the primitive sector,
$n(E)=1$, so the individual doubled family has either $(d,\ell)=(1,0)$
or $(0,1)$.
\end{proposition}

\begin{proof}
For one elementary transform with quotient line $G$ of degree $q$,
Grothendieck--Riemann--Roch and the triviality of the normal bundle of the
fiber give
\[
 \operatorname{ch}(i_*G)=[\Sigma_g]+q[\mathrm{pt}].
\]
Thus
\[
 c_1(E')=c_1(E)-[\Sigma_g],\qquad
 \operatorname{ch}_2(E')=\operatorname{ch}_2(E)-q[\mathrm{pt}].
\]
Since
\[
 p_1(\ad E)=4\operatorname{ch}_2(E)-c_1(E)^2,\qquad
 [\Sigma_g]^2=0,\qquad c_1(E)\cdot[\Sigma_g]=1,
\]
one obtains
\[
 p_1(\ad E')-p_1(\ad E)=2-4q=2(1-2q).
\]
All transforms are supported away from the framed outer circle, and hence
\begin{equation}\label{eq:relative-p1-langton}
              p_1(\ad F)=p_1(\ad E)+2\ell
\end{equation}
as a relative identity.  The fine-moduli property identifies
\[
 F\cong(\id_{\Sigma_g}\times f)^*\mathcal U\otimes p_D^*L_D.
\]
The adjoint removes the base-line factor.  After capping,
\eqref{eq:capped-degree} therefore says $p_1(\ad\widehat F)=-2d$.
By definition $p_1(\ad\widehat E)=-2n(E)$, so
\eqref{eq:relative-p1-langton} gives $n(E)=d+\ell$.

Two caps in the chosen chart are homotopic relative to the boundary.
A base-line twist leaves the adjoint bundle unchanged, and
$\ell=(p_1(\ad F)-p_1(\ad E))/2$ depends only on the framed endpoints.
These observations prove all independence assertions.  In the local
primitive sector, Corollary~\ref{cor:primitive-atom} gives $n(E)=1$.
Nonnegativity and integrality of $d$ and $\ell$ give the final dichotomy.
\end{proof}

\paragraph{The compactness distinction.}
The pair $(d,\ell)$ belongs to one individual doubled family and need not
be continuous under degeneration.  A sequence with $(d,\ell)=(1,0)$ can
converge in the quasimap compactification to $(0,1)$; the missing degree
then appears as a sphere component in the stable-map compactification.
That sphere and the quasimap basepoint are two records of the same
primitive charge and cannot be added as independent components.

Near rank drop, write a full-rank pencil in an affine base coordinate as
$s_\nu(z)=a_\nu z+b_\nu$.  If its limiting basepoint is
$z_\nu=x_\nu+iy_\nu$, choose a linear functional on the limiting image
line, absorb the phase of its transverse coefficient into $w_\nu$, and
write exactly
\begin{equation}\label{eq:rank-drop-exact-normal-form}
 s_\nu(z)=a_\nu(z-z_\nu)+\varepsilon_\nu w_\nu,
 \qquad \varepsilon_\nu>0,\quad \varepsilon_\nu\to0,\quad
 a_\nu\to v,\quad w_\nu\to w,\quad v\wedge w\ne0.
\end{equation}
After the seam-preserving change
$z=\varepsilon_\nu\zeta+iy_\nu$, the pencil is
\[
 [a_\nu(\zeta-\lambda_\nu)+w_\nu],
 \qquad \lambda_\nu=x_\nu/\varepsilon_\nu.
\]
If $\lambda_\nu$ has a finite limit, this converges to a full-rank mixed
bubble.  If $\lambda_\nu\to+\infty$ or $-\infty$, it is constant on
seam-centered compact sets and recentering at $z_\nu$ extracts,
respectively, a positive holomorphic sphere or a negative cylindrical
instanton.  Thus a seam basepoint is not a fourth type: center--scale
normalization either recovers a mixed bubble or one of the two pure ends.

The required analytic statement is strong convergence of arbitrary
primitive mixed sequences, in seam-compatible gauges and the first-order
norms used below, to this three-colored master compactification.  At
negative infinite seam ratio, affine recentering returns to the normalized
gauge rank-drop chart.  At positive infinite ratio the additional component
is an ordinary minimal sphere, whose nodal collar must be matched to the
finite-ratio charts.  Sections~\ref{sec:primitive-no-neck}--
\ref{sec:primitive-gluing} prove this analytic realization using
good-circle no-neck, the uniform anisotropic inverse, and local
surjectivity of the gluing maps.

\section{The algebraic primitive moduli}
\label{sec:primitive-algebraic}

\subsection{Mu\~noz's extension spaces}\label{sec:munoz-spaces}
Fix a line bundle $\Lambda$ of degree one on $\Sigma_g$.  Let
$J=J(\Sigma_g)$ be the Jacobian and let $\mathscr L$ be a Poincar\'e line
bundle on $\Sigma_g\times J$.  If $p:\Sigma_g\times J\to J$ is the
projection, put
\begin{equation}\label{eq:extension-bundle}
             \cE_g=R^1p_*(\mathscr L^2\otimes\Lambda^{-1}).
\end{equation}
With the standard principal polarization $\omega\in H^2(J;\ZZ)$, the
Grothendieck--Riemann--Roch calculation is
\begin{equation}\label{eq:ch-Eg}
                         \operatorname{ch}(\cE_g)=g+4\omega.
\end{equation}
The projective bundle
\[
                         N_g=\PP(\cE_g^\vee)\longrightarrow J
\]
parametrizes nonsplit extensions
\begin{equation}\label{eq:surface-extension}
 0\longrightarrow L\longrightarrow V
 \longrightarrow\Lambda\otimes L^{-1}\longrightarrow0.
\end{equation}
The middle term is stable and determines a holomorphic map
\begin{equation}\label{eq:i-map}
                              i_g:N_g\longrightarrow M_g.
\end{equation}

Mu\~noz studies stable rank-two bundles $E$ on
$\PP^1\times\Sigma_g$ with
\begin{equation}\label{eq:primitive-chern}
 c_1(E)=[\PP^1]+[\Sigma_g],\qquad c_2(E)=1,
 \qquad p_1(\ad E)=-2,
\end{equation}
for a polarization close to the $\Sigma_g$ boundary of the ample cone.  He
proves that their moduli space is the smooth projective bundle
\begin{equation}\label{eq:frakM}
 \frakM_g=\PP((\cE_g\oplus\cE_g)^\vee)\longrightarrow J,
 \qquad \dim_\CC\frakM_g=3g-1.
\end{equation}
See \cite[Sections~3--4]{Munoz99b}.  The two summands are the coefficients
of a section of $\cO_{\PP^1}(1)\otimes\cE_g$.

Fix $L\in J$, write $V=(\cE_g)_L$, and represent a point of the fiber by
\[
                      [v_0:v_1]\in\PP(V\oplus V).
\]
At $x=[x_0:x_1]\in\PP^1$, the corresponding surface extension has class
\begin{equation}\label{eq:extension-pencil}
                         v(x)=x_0v_0+x_1v_1.
\end{equation}

\begin{proposition}[rank stratification]\label{prop:rank-strata}
The moduli space $\frakM_g$ has the following two strata.
\begin{enumerate}
\item If $v_0,v_1$ are independent, \eqref{eq:extension-pencil} never
vanishes and defines a nonconstant parameterized minimal rational curve in
$M_g$.  Every parameterized minimal rational curve arises uniquely in
this way.
\item If $v_0,v_1$ are dependent, write
$(v_0,v_1)=(a_0v,a_1v)$.  The induced quasimap is constant away from the
single basepoint $[-a_1:a_0]$.  The rank-drop locus is the Segre bundle
\begin{equation}\label{eq:Segre}
 \cR_g=N_g\times_J\PP^1\longrightarrow\frakM_g
\end{equation}
of complex codimension $g-1$.
\end{enumerate}
For an individual primitive bundle, the two alternatives in
Proposition~\ref{prop:framed-splitting} correspond respectively to these
two strata.
\end{proposition}

\begin{proof}
If the two vectors are independent, no nonzero pair $(x_0,x_1)$ kills
their linear combination.  Thus \eqref{eq:extension-pencil} defines a
line in the projective extension fiber $\PP(V)$, followed by the map
$i_g$.  Mu\~noz's classification identifies this full-rank open locus with
the space of parameterized minimal curves.

If the two vectors are dependent, every nonzero extension away from
$[-a_1:a_0]$ represents the same point $[v]\in\PP(V)$.  At the unique
zero the surface bundle splits, and stable reduction records one
elementary modification.  Fiberwise, the dependent tensors form the
Segre image $\PP(V)\times\PP(\CC^2)$ in
$\PP(V\otimes\CC^2)$.  Its dimension is $(g-1)+1=g$, whereas the ambient
projective space has dimension $2g-1$, so its codimension is $g-1$.
The final assertion follows from the exact dichotomy in
Proposition~\ref{prop:framed-splitting}.
\end{proof}

For $g=1$, the extension space $V$ is one-dimensional, and every pair is
dependent.  The full-rank locus is empty and $\frakM_1=\cR_1$.  This is
the algebraic reason the genus-one primitive class is invisible in the
point-valued character variety.

\subsection{Evaluation at infinity and the dimension ledger}
\label{sec:evaluation-infinity}
Choose $\infty=[1:0]\in\PP^1$.  On the locus $v_0\neq0$, evaluation of
the induced stable map or quasimap is
\begin{equation}\label{eq:ev-infinity}
 \ev_\infty([v_0:v_1])=i_g([v_0]).
\end{equation}
Hence every primitive bubble, full rank or rank drop away from infinity,
attaches through $i_g(N_g)$.

Fix $[v]\in N_g$.  After the overall projective scalar is used to set
$v_0=v$, the full-rank fiber of $\ev_\infty$ is
\[
                              V\setminus\CC v.
\]
The real affine group preserving the two half-planes and fixing infinity
is
\begin{equation}\label{eq:GR}
 G_\RR=\{w\longmapsto aw+ib:a>0,\ b\in\RR\}.
\end{equation}
It fixes the translation and scale of a generic mixed bubble.  The reduced
fiber therefore has real dimension $2g-2$.

The dimensions of the attachment spaces are
\begin{equation}\label{eq:dimensions-N-M}
             \dim_\CC N_g=2g-1,
       \qquad\dim_\CC M_g=3g-3.
\end{equation}
Mu\~noz proves that $i_g$ is generically finite onto its image.  Thus
\begin{equation}\label{eq:image-codim}
                 \codim_\RR i_g(N_g)=2(g-2).
\end{equation}
Let $e_0:\RR\to M_g$ be the matching evaluation path of a residual
index-zero mixed solution.  Whenever it meets the regular part of
$i_g(N_g)$, the reduced full-rank primitive stratum has dimension
\begin{equation}\label{eq:bubble-dim-ledger}
                 (2g-2)+1-2(g-2)=3.
\end{equation}
The interpolation parameter adds one dimension and the integral $p_1$
incidence removes four, leaving isolated boundary points.

For $g\geq3$, however, a generic one-dimensional path misses the
positive-codimension image in \eqref{eq:image-codim}.  The matching
correction is therefore absent.  In genus two, $N_2$ and $M_2$ have the
same dimension and $i_2$ has a nonzero degree, computed next.

\subsection{The genus-two degree}\label{sec:genus-two-degree}
Let $h=c_1(\cO_{N_2}(1))$, and normalize the principal polarization by
\begin{equation}\label{eq:omega-normalization}
                    \int_{J(\Sigma_2)}\frac{\omega^2}{2}=1.
\end{equation}
If $\alpha\in H^2(M_2;\ZZ)$ is the standard generator, Mu\~noz's
universal-bundle formulas give
\begin{equation}\label{eq:munoz-pullback}
 i_2^*\alpha=4\omega+h,
 \qquad h^2+4\omega h+8\omega^2=0.
\end{equation}

\begin{proposition}\label{prop:degree-four}
The map $i_2:N_2\to M_2$ has degree four.
\end{proposition}

\begin{proof}
Using \eqref{eq:munoz-pullback} and $\omega^3=0$ gives
\[
 h^3=8\omega^2h,
 \qquad 12\omega h^2=-48\omega^2h.
\]
Therefore
\begin{equation}\label{eq:N2-intersection}
 \int_{N_2}(4\omega+h)^3
 =\int_{N_2}8\omega^2h=16,
\end{equation}
where the last equality uses \eqref{eq:omega-normalization} and integration
over the projective-line fiber.  The genus-two moduli space $M_2$ is a
smooth intersection of two quadrics in $\PP^5$, and $\alpha$ is its
hyperplane class.  Hence
\[
                         \int_{M_2}\alpha^3=4.
\]
The ratio of \eqref{eq:N2-intersection} to this number is the degree:
$\deg i_2=16/4=4$.
\end{proof}

All local algebraic multiplicities are positive in the complex
orientation.  Thus, after the analytic incidence determinant is
identified with its algebraic counterpart, the magnitude of the integral
genus-two contribution is four.

\subsection{The genus-one excess line}\label{sec:genus-one-excess}
For $g=1$, \eqref{eq:ch-Eg} reads
\begin{equation}\label{eq:torus-E}
                         \operatorname{ch}(\cE_1)=1+4\omega.
\end{equation}
The two coefficients of the primitive extension form the rank-two bundle
\begin{equation}\label{eq:Ezeta}
 \cE_\zeta=H^0(\PP^1,\cO(1))\otimes\cE_1
           \cong\cE_1\oplus\cE_1.
\end{equation}
Consequently
\begin{equation}\label{eq:torus-degree-eight}
 c_1(\det\cE_\zeta)=8\omega,
 \qquad
 \int_{J(\Sigma_1)}c_1(\det\cE_\zeta)=8.
\end{equation}
After the rank-drop basepoint and the real affine symmetry are fixed, the
elliptic Jacobian is the excess primitive family.  Section
\ref{sec:determinant-locality} identifies its normal incidence determinant
with $\det\cE_\zeta$.  Thus \eqref{eq:torus-degree-eight} is the magnitude
of the integral genus-one contribution.

\subsection{Cylinder calibration}\label{sec:cylinder-calibration}
The algebraic magnitudes can be calibrated, including sign, without
assuming the point theorem.  Let $a_G,b_G$ denote Mu\~noz's surface and
point generators in the instanton Floer ring of $\Sigma_g\times S^1$, and
let $a_S,b_S$ be the geometric generators of $QH^*(M_g)$.  These act on
the gauge and symplectic Floer groups of every admissible target containing
the chosen surface component.

We record the normalization dictionary.  Mu\~noz uses
$\mu(q)=-\frac14p_1/[q]$ and defines his point generator by
$b=-4\mu(q)=p_1/[q]$; his surface generator is
$a=2\mu(\Sigma)=-\frac12p_1/[\Sigma]$
\cite[equation~(4)]{Munoz99a}.  Thus $b_G,b_S$ are exactly the positive
integral point generators used here.  DFL's comparison for
$p_1/[\Sigma]=-2a$ implies, over $\QQ$,
\begin{equation}\label{eq:surface-compatibility}
                              Na_G=a_SN.
\end{equation}
For genus two, Mu\~noz's corrected quantum generator is
$\widehat b_S=b_S+4$, and his relation
$a_S^2+\widehat b_S-8=0$ is therefore
$a_S^2+b_S-4=0$ in the unshifted geometric generator
\cite[Example~24]{Munoz99b}.  On the gauge side his corresponding relation
is $a_G^2+b_G-8=0$ \cite[Example~22]{Munoz99a}.  In genus one his gauge
relation is $b_G=8\id$ \cite[Lemma~11]{Munoz99a}, whereas the geometric
class $b_S$ vanishes because $M_1$ is a point
\cite[Example~23]{Munoz99b}.

\begin{proposition}[low-genus cylinder identities]
\label{prop:cylinder-identities}
In the positive integral $p_1$ normalization,
\begin{equation}\label{eq:cylinder-identities}
 Nb_G-b_SN=
 \begin{cases}
 8N,&g=1,\\
 4N,&g=2.
 \end{cases}
\end{equation}
\end{proposition}

\begin{proof}
For $g=1$, the odd-determinant character variety is a point, so $b_S=0$.
Mu\~noz's cylinder presentation has $b_G=8\id$, proving the first formula.

For $g=2$, the instanton and quantum relations in the positive integral
normalization are
\[
                    a_G^2+b_G-8=0,
              \qquad a_S^2+b_S-4=0.
\]
Acting on the two modules and using \eqref{eq:surface-compatibility}
twice gives
\[
 \begin{aligned}
 Nb_G-b_SN
 &=N(8-a_G^2)-(4-a_S^2)N\\
 &=4N-a_S^2N+Na_G^2=4N.
 \end{aligned}
\]
\end{proof}

The proposition already proves that uncorrected point compatibility is
false in low genus.  In the final orientation argument it has a second,
noncircular role: determinant excision makes the local bubble coefficient
universal, and \eqref{eq:cylinder-identities} fixes the universal DFL sign.

\section{Dirichlet reconstruction and primitive gluing}
\label{sec:reconstruction}

\subsection{Seam-admissible product bundles}\label{sec:seam-admissible}
Decompose the base sphere into closed hemispheres
\begin{equation}\label{eq:hemispheres}
                         \PP^1=H_-\cup_\Gamma H_+,
\end{equation}
where $\Gamma$ is the compactified matching line and
$\infty\in\Gamma$ is the attaching point.  The gauge equation is placed
on $H_-\times\Sigma_g$ and the holomorphic-map equation on $H_+$.  A
primitive bundle $E\in\frakM_g$ induces on the stable fibers of the base a
classifying map to $M_g$.

\begin{definition}\label{def:seam-admissible}
A bundle $E\in\frakM_g$ is \emph{seam admissible} if every restriction
$E|_{\{z\}\times\Sigma_g}$ with $z\in\Gamma$ is stable and the
classifying quasimap has no basepoint in $H_+$.  We denote the
seam-admissible locus by $\frakM_g^{\rm sa}$.
\end{definition}

Every full-rank pair is seam admissible.  A rank-drop pair is seam
admissible precisely when its unique basepoint lies in the interior of
$H_-$.  A basepoint in $H_+$ makes the positive map undefined, while a
basepoint on $\Gamma$ cannot occur in a smooth seam-admissible object.  For
a sequence approaching $\Gamma$, the center--scale normal form
\eqref{eq:rank-drop-exact-normal-form} either recovers a full-rank mixed
bubble at finite seam ratio or converges to one of the two pure limits.
The real M\"obius group preserving
\eqref{eq:hemispheres} acts transitively on the interior of $H_-$; after
fixing infinity, its affine subgroup $G_\RR$ in \eqref{eq:GR} normalizes a
gauge-side basepoint to a compact set.

\subsection{Exact reconstruction}\label{sec:exact-reconstruction}
For $\rho>0$, equip the gauge hemisphere with the product K\"ahler form
\begin{equation}\label{eq:omega-rho}
                    \omega_\rho=\rho^2\omega_{H_-}+\omega_{\Sigma_g}.
\end{equation}
This is the form obtained by pulling the original metric back under base
dilation by $\rho$.

\begin{theorem}[Dirichlet reconstruction]\label{thm:reconstruction}
Let $E\in\frakM_g^{\rm sa}$.  For every $\rho>0$ there is a unique
fixed-determinant projectively ASD connection $A_{E,\rho}$ on
$E|_{H_-\times\Sigma_g}$ whose boundary metric is the family of
Narasimhan--Seshadri metrics on the stable seam fibers.  Together with the
holomorphic classifying map
\[
                         u_E:H_+\longrightarrow M_g,
\]
it is an exact solution of the $\rho$-rescaled mixed equation.  It has
relative clutching charge one and normalized mixed energy $\frac12$.  For
$\rho>0$ the solution depends smoothly on $(E,\rho)$.
\end{theorem}

\begin{proof}
Each stable odd-degree bundle on a seam fiber has a unique
fixed-determinant Hermitian--Einstein metric.  Simplicity removes every
noncentral ambiguity, and elliptic parameter dependence makes these
metrics a smooth family $H_\Gamma$ on
$E|_{\Gamma\times\Sigma_g}$.

Donaldson's boundary-value theorem
\cite[Theorem~1]{Donaldson92} gives a unique Hermitian metric
$H_{E,\rho}$ with fixed determinant satisfying
\begin{equation}\label{eq:HYM-Dirichlet}
 H_{E,\rho}|_{\Gamma\times\Sigma_g}=H_\Gamma,
 \qquad
 i\Lambda_{\omega_\rho}F_{H_{E,\rho}}^\perp=0.
\end{equation}
On a K\"ahler surface, integrability and the trace-free Hermitian
Yang--Mills equation are equivalent to projective anti-self-duality.  The
Chern connection of $H_{E,\rho}$ is therefore the required
$A_{E,\rho}$.  No stability of $E$ with respect to $\omega_\rho$ is
needed: the Dirichlet theorem applies to every holomorphic bundle on a
compact K\"ahler manifold with boundary.  This point is essential because
the polarization \eqref{eq:omega-rho} approaches the chamber opposite to
the one used to define $\frakM_g$.

On $H_+$, seam admissibility gives the holomorphic map $u_E$.  Its
pullback universal bundle is $E|_{H_+\times\Sigma_g}$, and its canonical
flat representative along $\Gamma$ is the Chern connection of
$H_\Gamma$.  Thus the two sides satisfy the exact matching condition.
Holomorphicity gives the positive equation.

The Chern--Weil filling identity and $p_1(\ad E)=-2$ give relative charge
one and energy $\frac12$.  The linearization of
\eqref{eq:HYM-Dirichlet} in the Hermitian metric direction is the
Dirichlet Chern Laplacian on trace-free self-adjoint endomorphisms.  Its
kernel is zero by integration by parts.  Elliptic parameter dependence
and the implicit-function theorem give smooth dependence on $(E,\rho)$.
\end{proof}

For a full-rank $E$, the positive map is part of a basepoint-free minimal line and
the negative half is its unique Dirichlet ASD completion.  For a rank-drop
$E$ with basepoint in $H_-$, the positive map is constant and the entire
modification lies on the gauge hemisphere.  Thus every individual
algebraic type in Proposition~\ref{prop:framed-splitting} is realized by
an exact local mixed
solution.  In genus one the second type is the pure gauge bubble.

\subsection{Fredholm reduction}\label{sec:fredholm-reduction}
Fix $E\in\frakM_g^{\rm sa}$ and $\rho>0$.  Infinitesimal deformations of
the reconstructed pair split into a Dolbeault deformation of $E$ and a
trace-free Hermitian-metric deformation with zero Dirichlet value.  The
metric block is invertible by the last paragraph of the proof of
Theorem~\ref{thm:reconstruction}.  Domain reparameterizations are added
only after this fixed-domain problem is coupled to the residual solution.

\begin{proposition}[fixed-scale Fredholm-to-Dolbeault reduction]
\label{prop:fredholm-dolbeault}
Let $\cD^{\rm mix}_{E,\rho}$ be DFL's gauge-fixed mixed operator, with
the relative boundary condition and canonical linearized matching
correspondence in its domain.  Infinitesimal holomorphic doubling gives
a natural real-linear isomorphism
\begin{equation}\label{eq:dolbeault-tangent-obstruction}
 \ker\cD^{\rm mix}_{E,\rho}
 \cong H^1(\PP^1\times\Sigma_g,\End_0E).
\end{equation}
The mixed operator is surjective.  Hence the fixed-domain mixed bubble is
regular and its determinant line has the complex orientation transported
through \eqref{eq:dolbeault-tangent-obstruction}.  Evaluation at infinity and
quotienting by $G_\RR$ are separate finite-dimensional operations imposed
only after coupling to the residual component.
\end{proposition}

\begin{proof}
DFL's operator has components
\[
                 (d_A^+\zeta,-d_A^*\zeta,D_u\nu).
\]
Its domain satisfies
$*\zeta|_\Gamma=0$ and
$(\zeta|_{\{z\}\times\Sigma_g},\nu(z))\in T\cL_{\rm NS}$, while its
formal adjoint uses the same canonical correspondence
\cite[(1.11), (1.13), and (5.1)]{DaeFukLip21b}.  Their local estimates,
patched over the compact hemispheres, make this a Fredholm boundary
problem.

DFL regularity makes every kernel element smooth.  The K\"ahler identity
identifies the negative component $a^{0,1}$ with a Dolbeault cocycle.  On
the positive side, harmonic inclusion for the universal connection
identifies a map variation with the infinitesimal deformation of the
pulled-back universal bundle.  The matching condition says that the two
seam classes agree; their difference is
$\bar\partial_\alpha\xi$ for a smooth infinitesimal complex gauge
transformation.  Extend $\xi$ smoothly across a collar and differentiate
the finite-scale doubling construction of Lemma~\ref{lem:doubling}.
This produces the class in the right side of
\eqref{eq:dolbeault-tangent-obstruction}.  Different extensions change it
by a Dolbeault coboundary.

If this class vanishes, the positive harmonic variation is zero and the
negative variation changes only the Hermitian metric of a fixed
holomorphic bundle.  Its boundary value is zero.  The linearized
Dirichlet HYM equation has
\[
 \langle L^{\rm Dir}_{E,\rho}s,s\rangle
       =2\|\bar\partial_{A_{E,\rho}}s\|^2.
\]
It has zero kernel, proving injectivity.  Conversely, Mu\~noz identifies
the primitive moduli with the smooth projective bundle
$\PP(\cE_\zeta^\vee)\to J(\Sigma_g)$.  Integrate a class in $H^1$ to a
curve in this projective bundle, apply Theorem~\ref{thm:reconstruction}
at the fixed scale, and differentiate.  This proves surjectivity in
\eqref{eq:dolbeault-tangent-obstruction}.

It remains to rule out a DFL cokernel.  Deform the coefficients to product
form in a seam collar and use DFL's bounded domain-straightening maps
\cite[Proposition~5.27]{DaeFukLip21b}.  The K\"ahler identities identify
the gauge symbol and Green form with the Dolbeault--Dirac symbol and Green
form.  Fiberwise Hodge decomposition pairs every nonzero exact vertical
mode with its coexact mate; the two half-line operators are adjoints and
have cancelling indices.  Irreducibility removes the zero exact mode.
The remaining harmonic matching condition is the graph transmission
condition.  Its Sobolev domain is the domain of the uncut closed operator
\cite[Example~3.23]{BaerBallmann16}.  Relative index gluing therefore gives
\[
 \ind_\RR\cD^{\rm mix}_{E,\rho}
 =\ind_\RR\sqrt2(\bar\partial_E^*\oplus\bar\partial_E)
 =2(h^1-h^0-h^2).
\]
Stability gives $h^0=0$, while Riemann--Roch gives
\[
 h^1-h^2=3g-1.
\]
The smooth projective-bundle description gives $h^1=3g-1$, including on
the Segre locus.  Thus the real kernel dimension equals the real index,
and the cokernel is zero.  DFL's adjoint regularity and its identical
Lagrangian matching condition show that this is the analytic cokernel.
Since the cokernel vanishes, \eqref{eq:dolbeault-tangent-obstruction}
transports the complex orientation to the determinant line.
\end{proof}

The rank-drop normal can be read explicitly.  At a tensor
$[v\otimes a]\in\PP(V\otimes\CC^2)$, let
$\ell_v=\CC v$ and $\ell_a=\CC a$.  The normal to the Segre locus is
\begin{equation}\label{eq:Segre-normal}
 \Hom(\ell_v,V/\ell_v)\otimes
 \Hom(\ell_a,\CC^2/\ell_a),
\end{equation}
of complex dimension $g-1$.  In genus two it is a complex line; after its
basepoint is fixed, both real components remain as smoothing directions
from rank drop to a full-rank bubble.  The seam-preserving affine group has
no rotational factor and does not remove the phase.  Thus the genus-two
Segre locus has real codimension two inside the primitive face.  In genus
one the normal has rank zero, leaving the elliptic excess family.

\subsection{Primitive no-neck at good circles}\label{sec:primitive-no-neck}
Let $(A_\nu,u_\nu)$ be a sequence with one primitive matching center
$z_\nu\to z_0$ and concentration scales $\rho_\nu\to0$, and suppose its
doubled bundles have a limit $E$ in the three-colored master
compactification described after Proposition~\ref{prop:framed-splitting}.
Write
$\cE_\nu(a,b)$ for the normalized mixed energy in the paired half-annulus
$a<|z-z_\nu|<b$.

\begin{proposition}[primitive good-circle no-neck]
\label{prop:primitive-no-neck}
After passage to a subsequence, there are radii
\[
                 \rho_\nu\ll a_\nu\ll b_\nu\longrightarrow0
\]
such that
\begin{equation}\label{eq:no-neck}
                         \cE_\nu(a_\nu,b_\nu)\longrightarrow0.
\end{equation}
The capped region inside $a_\nu$ has clutching sector one and energy
tending to $\frac12$.  Hence no positive-energy atom or further analytic
scale lies between the primitive core and the residual solution.  On the
two cut circles, the spliced representatives converge in a common seam
gauge to the same flat attaching representative.
\end{proposition}

\begin{proof}
The stable limit is regular at infinity.  For a sufficiently large
rescaled radius $R$, cap its tail in a contractible fine-moduli chart.
The finite charge identity in Proposition~\ref{prop:framed-splitting}
shows that, for $R$ and then $\nu$
large, the capped region $|z-z_\nu|\leq R\rho_\nu$ has clutching sector
one.  This includes rank drop because the lost map degree is exactly the
length-one elementary modification.  The local charge theorem gives the
same sector for the outer atom.  Every annulus between the two caps
therefore has sector zero.

Put $E_\nu(r)=\cE_\nu(0,r)$.  First choose the outer cuts from smooth
convergence on compact subsets of the punctured disc to the residual
solution.  DFL removability extends that solution over $z_0$, so a
diagonal choice gives $b_\nu\to0$, with
$b_\nu/\rho_\nu\to\infty$, such that the outer-tail energy relative to
the residual solution and $b_\nu E_\nu'(b_\nu)$ both tend to zero.

Since
\[
 \int_{r_1}^{r_2}rE_\nu'(r)\,d\log r
       =E_\nu(r_2)-E_\nu(r_1),
\]
choose $R_\nu\to\infty$ slowly enough that
$R_\nu^2\rho_\nu<b_\nu$, and subdivide the logarithmic region from
$R_\nu\rho_\nu$ to $R_\nu^2\rho_\nu$ into dyadic shells.  Bounded total
energy and coarea give a shell of energy tending to zero and a radius
$a_\nu$ in it for which
\begin{equation}\label{eq:good-circle}
 a_\nu E_\nu'(a_\nu)+b_\nu E_\nu'(b_\nu)\longrightarrow0.
\end{equation}

At either good circle form DFL's spliced connection $B_{\nu,r}$ on
$S^1\times\Sigma_g$: it joins the gauge semicircle to a connection
representing the holomorphic semicircle through the flat interpolation
strip.  The small-energy mean-value estimate on the selected shell and
the proof of \cite[Lemma~4.62]{DaeFukLip21b} give the localized slice bound
\[
 |\CS_{\beta_{\nu,r}}(B_{\nu,r})|
             \leq C rE_\nu'(r).
\]
The flat references at the two radii are joined by the seam path, so the
Chern--Simons functional is independent of that choice.  The Stokes
calculation of \cite[Lemma~4.67]{DaeFukLip21b}, now on the annulus, gives
\[
 \cE_\nu(a_\nu,b_\nu)
 =\CS(B_{\nu,b_\nu})-\CS(B_{\nu,a_\nu})
\]
with the normalization of \eqref{eq:top-energy}.  The possible integral
change of lift is zero because the annulus has clutching sector zero.
Equation~\eqref{eq:good-circle} proves \eqref{eq:no-neck}.

The total local energy tends to $\frac12$, the annular energy vanishes,
and the outer-tail energy tends to that of the removable residual
solution.  Thus the sector-one inner core has energy $\frac12$ in the
limit.  A further atom would have positive
half-integral charge, whereas a diffuse remainder is excluded by the
same Chern--Simons identity.  Finally, DFL's small-energy Coulomb gauges
on the two selected shells converge to the flat seam representatives;
stable reduction at infinity identifies both with the attaching value.
\end{proof}

\subsection{Inverse-adiabatic Dirichlet estimates}
\label{sec:uniform-inverse}
Put $\tau=\rho^2$.  Since the K\"ahler form is a product, the trace-free
HYM equation is equivalent to
\begin{equation}\label{eq:tau-HYM}
 i\Lambda_{H_-}F_H^\perp
       +\tau i\Lambda_{\Sigma_g}F_H^\perp=0.
\end{equation}
Thus the limit $\tau=0$ is elliptic in the base variable and regards the
surface variable as a parameter.  This is the useful direction of the
inverse-adiabatic limit.

For the following statements, $W^{r,s;p}_{B,\Sigma}$ controls mixed
derivatives with at most $r$ base derivatives and at most $s$ surface
derivatives in fixed product charts.  On the HYM-equation
component we use the equivalent $\rho$-weighted target norm obtained by
multiplying the equation by $\rho^2$, as in
\eqref{eq:tau-HYM}.  Fixed weights are added at the attaching end.

\begin{theorem}[inverse-adiabatic Dirichlet regularity]
\label{thm:inverse-adiabatic}
Let $C\subset\frakM_g^{\rm sa}/G_\RR$ be compact.  The Dirichlet metrics
$H_{E,\rho}$ extend smoothly in $(E,\tau)$ to $\tau=0$ in every
anisotropic $W^{r,s;p}_{B,\Sigma}$ space.  The limiting metric $H_{E,0}$ is
the unique fixed-determinant metric satisfying
\begin{equation}\label{eq:base-flat-limit}
 i\Lambda_{H_-}F_{H_{E,0}}^\perp=0,
 \qquad H_{E,0}|_{\Gamma\times\Sigma_g}=H_\Gamma.
\end{equation}
In particular, relative to a smooth reference family, the metrics,
connections, and all fixed-coordinate derivatives are bounded uniformly
for $E\in C$ and $0\leq\rho\leq\rho_0$.  This conclusion holds through
the rank-drop locus.
\end{theorem}

\begin{proof}
Fix $x\in\Sigma_g$.  The restriction of $E$ to
$H_-\times\{x\}$ is a holomorphic bundle on a disc.  Donaldson's
Dirichlet theorem in complex dimension one gives the unique
fixed-determinant solution of \eqref{eq:base-flat-limit} with the
prescribed boundary metric.  Equivalently, it is the harmonic metric for
a disc with target the nonpositively curved symmetric space of positive
Hermitian matrices of determinant one.  Ordinary parameter-dependent
Dirichlet theory makes these solutions smooth in $(E,x)$ and gives a
smooth metric $H_{E,0}$ on the product.

Write nearby metrics as $H_{E,0}e^s$, where $s$ is trace-free,
self-adjoint, and zero on the boundary.  The derivative at
$(s,\tau)=(0,0)$ of the left side of \eqref{eq:tau-HYM} is the
base Dirichlet Chern Laplacian
\[
              L^B_{E,0}s
       =i\Lambda_{H_-}\bar\partial_B\partial_{H_{E,0},B}s.
\]
For each $x$, integration by parts identifies its quadratic form with
$2\|\bar\partial_Bs\|^2$.  Its kernel is zero by the boundary condition,
and the Dirichlet operator on the disc has index zero.  It is therefore
an isomorphism.  Differentiating in the parameter $x$ gives the standard
parameter-elliptic estimates, so
\[
 L^B_{E,0}:W^{r+2,s;p}_{B,\Sigma,0}
       \longrightarrow W^{r,s;p}_{B,\Sigma}
\]
is an isomorphism for every $r,s$ and $p>1$.

It remains to justify uniformity for the exact solutions rather than use
an implicit-function theorem with a surface-derivative loss.  Put
$h_\rho=H_{E,0}^{-1}H_{E,\rho}$.  Because $H_{E,0}$ is base-flat,
$\Lambda_{\omega_\rho}F_{H_{E,0}}^\perp$ is bounded independently of
$\rho$.  The Chern-connection trace inequality gives
\[
 -\Delta_{\omega_\rho}\log\tr h_\rho\leq C.
\]
If $-\Delta_{H_-}\phi=1$ and $\phi|_\Gamma=0$, then
$-\Delta_{\omega_\rho}(\rho^2\phi)=1$.  Since $h_\rho=\id$ on the
boundary, the maximum principle gives
\begin{equation}\label{eq:metric-C0-close}
                 \log\tr h_\rho\leq\log2+C\rho^2.
\end{equation}
The determinants agree, so \eqref{eq:metric-C0-close} makes
$s_\rho=\log h_\rho$ uniformly small in $C^0$.

This initial smallness gives the sharper order in $\tau$.  Define
\[
 \Phi(H,\tau)=i\Lambda_{H_-}F_H^\perp
              +\tau i\Lambda_{\Sigma_g}F_H^\perp.
\]
The fundamental theorem of calculus along
$H_{E,0}e^{t s_\rho}$ gives the exact equation
\begin{equation}\label{eq:integrated-HYM-linearization}
 \left(\int_0^1
   D_H\Phi(H_{E,0}e^{t s_\rho},\tau)\,dt\right)s_\rho
 =-\tau i\Lambda_{\Sigma_g}F_{H_{E,0}}^\perp.
\end{equation}
For $\rho$ small, the averaged operator has the same uniform base
Dirichlet coercivity as $L^B_{E,0}$; the $C^0$ smallness permits its
coefficient error to be absorbed.  Base Dirichlet regularity therefore
first gives $\|s_\rho\|_{W^{2,0;p}_{B,\Sigma}}\leq C\tau$.

Apply surface and parameter difference quotients to
\eqref{eq:integrated-HYM-linearization}.  At every order the highest
derivative $v$ satisfies
the linearized Dirichlet equation, whose normalized quadratic form obeys
\begin{equation}\label{eq:anisotropic-energy}
 \|\nabla_Bv\|_{L^2}^2
       +\tau\|\nabla_\Sigma v\|_{L^2}^2
 \leq C\bigl(\|f\|_{H^{-1}_B}^2+\|v\|_{L^2}^2\bigr).
\end{equation}
Here $f$ contains only lower-order difference quotients.  The base
Dirichlet Poincar\'e inequality absorbs the last term, with a constant
independent of $\tau$.  Difference-quotient induction, followed by base
Dirichlet regularity and Sobolev embedding, gives
\[
 \|s_\rho\|_{W^{r,s;p}_{B,\Sigma}}\leq C_{r,s,p}\tau
\]
at every finite pair of orders.  At each induction step, derivatives of
the averaged operator multiply already controlled lower derivatives of
$s_\rho$, while the differentiated right side of
\eqref{eq:integrated-HYM-linearization} retains its factor $\tau$.
Differentiating the equation in $\tau$ and repeating the same estimate
gives all $\tau$ derivatives.  Compactness of $C$ makes all constants
uniform.  This proves the claimed smooth extension, and
Donaldson uniqueness identifies it with the family in
Theorem~\ref{thm:reconstruction}.  Rank drop
creates an unstable \emph{surface} fiber, but the limiting Dirichlet
operator acts in the base disc, where every restricted bundle is
holomorphically trivial; hence it introduces no exceptional model.
\end{proof}

\begin{proposition}[uniform metric-block coercivity]
\label{prop:metric-coercivity}
Let $L_{E,\rho}^{\rm Dir}$ be the trace-free self-adjoint Dirichlet
Chern Laplacian at a reconstructed bubble.  There is a constant $c>0$,
independent of $E\in\frakM_g^{\rm sa}$ and $0<\rho\leq1$, such that
\begin{equation}\label{eq:metric-coercivity}
 \langle L_{E,\rho}^{\rm Dir}s,s\rangle
 \geq c\rho^{-2}\|s\|_{L^2(\omega_\rho)}^2,
 \qquad s|_{\Gamma\times\Sigma_g}=0.
\end{equation}
Consequently its $L^2$ inverse has norm at most $c^{-1}\rho^2$.
\end{proposition}

\begin{proof}
The Bochner--Kodaira identity on endomorphisms gives
\[
 2\|\bar\partial_{A_{E,\rho}}s\|^2
 =\|\nabla_{A_{E,\rho}}s\|^2
  +\langle i[\Lambda_{\omega_\rho}F_{A_{E,\rho}},s],s\rangle.
\]
The last term vanishes because the connection is projectively HYM.
Kato's inequality and the scalar Dirichlet Poincar\'e inequality on
$H_-$ imply
\[
 \|\nabla_{A_{E,\rho}}s\|_{L^2(\omega_\rho)}^2
 \geq \|d|s|\|_{L^2(\omega_\rho)}^2
 \geq c\rho^{-2}\|s\|_{L^2(\omega_\rho)}^2.
\]
This argument contains no curvature bound and is therefore unaffected by
rank drop.
\end{proof}

\begin{theorem}[uniform reduced inverse]\label{lem:uniform-inverse}
Let $C$ be a compact subset of $\frakM_g^{\rm sa}/G_\RR$.  On the
complement of the Dolbeault tangent space, the mixed linearization at the
reconstructed bubbles has a right inverse $Q_{E,\rho}$ satisfying
\begin{equation}\label{eq:uniform-inverse}
                         \|Q_{E,\rho}\|\leq C_0
\end{equation}
for $E\in C$ and $0<\rho\leq\rho_0$.
\end{theorem}

\begin{proof}
Fix $p>4$ and put $\tau=\rho^2$.  We first specify the norms in the
statement.  After the collar Cauchy--Riemann identification below, write
$s$ for the full coexact--metric pair and $\chi$ for the full
exact--unitary-gauge pair; each displayed scalar norm is applied
componentwise.  The first pair, with zero Dirichlet value, has norm
\[
 \begin{split}
 \|s\|_{\mathcal M_\tau^p}={}&\|s\|_{L^p}
 +\|\nabla_Bs\|_{L^p}+\|\nabla_B^2s\|_{L^p}
 +\rho\|\nabla_\Sigma s\|_{L^p}\\
 &+\rho\|\nabla_B\nabla_\Sigma s\|_{L^p}
 +\tau\|\nabla_\Sigma^2s\|_{L^p}.
 \end{split}
\]
For the exact--gauge pair, with covariant Neumann boundary
condition, use
\[
 \|\chi\|_{\mathcal G_\tau^p}
 =\|\chi\|_{W^{2,p}_\Sigma}
 +\rho^{-1}\|\nabla_B\chi\|_{W^{1,p}_\Sigma}
 +\rho^{-2}\|\nabla_B^2\chi\|_{L^p}.
\]
The Dolbeault component has its fixed $W^{1,p}$ graph norm.  On the target
leave the integrability, positive Cauchy--Riemann, and Coulomb rows
unscaled, and multiply only the scalar HYM row by $\rho^2$.  These direct
sum norms are denoted by $\mathcal X_{E,\tau}^p$ and
$\mathcal Y_\tau^p$.

We next verify that these are norms on DFL's first-order boundary problem,
rather than on a formal second-order replacement.  In DFL's seam collar,
write a spatial one-form as $b=q+\tau_0ds$, where $\tau_0$ is a field
component.  Their Hodge decomposition (5.17)--(5.21) is
\begin{equation}\label{eq:dfl-collar-hodge}
 q=h+d_\alpha\zeta+*_\Sigma d_\alpha\zeta'.
\end{equation}
Here $h$ is harmonic and
\[
 \zeta,\zeta'\in
 W^{1,p}_{s,\theta}W^{1,p}_\Sigma
 \cap L^p_{s,\theta}W^{2,p}_\Sigma.
\]
The boundary conditions are
\begin{equation}\label{eq:dfl-collar-boundary}
 \tau_0|_{s=0}=0,\qquad \zeta'|_{s=0}=0,\qquad
 (h|_{s=0},\nu|_{s=0})\in V_\alpha,
\end{equation}
where $V_\alpha$ is the harmonic Narasimhan--Seshadri graph.  The exact
potential $\zeta$ is unrestricted at the seam.  The surface zero-form
Laplacian has a uniform positive gap because every seam connection is
irreducible.  Hence \eqref{eq:dfl-collar-hodge} and its inverse are
uniformly bounded in the displayed anisotropic spaces.  This works with
the full collar fields and requires no extension of a seam first jet.

Freeze the coefficients and take
$d_\alpha^*d_\alpha e=\mu e$, $\mu>0$.  Substitution of
\[
 q=x\mu^{-1/2}d_\alpha e
   +y\mu^{-1/2}*_\Sigma d_\alpha e
\]
in DFL (5.7) pairs $(y,\tau_0)$ and $(x,\varphi)$ into two
Cauchy--Riemann matrices.  After the $\theta$ derivative is added and the
system is squared, the off-diagonal first derivatives cancel.  With the
row normalization above, the two scalar symbols are
\begin{equation}\label{eq:two-parameter-symbols}
 |\xi_B|^2+\rho^2\mu,\qquad
 \rho^{-2}|\xi_B|^2+\mu.
\end{equation}
The first has the Dirichlet condition in
\eqref{eq:dfl-collar-boundary}; the second has the complementary Neumann
condition supplied by the adjoint equation.  They are respectively the
metric and unitary-gauge blocks.  Base Dirichlet Poincar\'e gives a gap for
the first symbol, uniformly when $\rho^2\mu$ tends to zero.  The lower bound
$\mu\geq\mu_0>0$ gives a gap for the second symbol, including its
base-constant mode.  Reflection and the multipliers in
\eqref{eq:two-parameter-symbols} give
\begin{equation}\label{eq:two-block-estimates}
 \|s\|_{\mathcal M_\tau^p}
 \leq C\|L_{E,\tau}^{\rm Dir}s\|_{L^p},
 \qquad
 \|\chi\|_{\mathcal G_\tau^p}
 \leq C\|G_{E,\tau}^{\rm Neu}\chi\|_{L^p}.
\end{equation}

The harmonic term in \eqref{eq:dfl-collar-hodge}, coupled to $\nu$ by
$V_\alpha$, is the Dolbeault transmission block of
Proposition~\ref{prop:fredholm-dolbeault}.  On the complement of its kernel,
its smallest singular value is uniformly positive over $C$.  DFL's
domain-straightening map in Proposition 5.27 acts separately on the
harmonic, exact, and coexact entries of the full collar field.  Its formula
and inverse are uniformly bounded in the norms above.  A partition of
unity, coefficient freezing, and \eqref{eq:two-block-estimates} therefore
give
\begin{equation}\label{eq:uniform-dfl-with-compact-term}
 \|x\|_{\mathcal X_{E,\tau}^p}
 \leq C\bigl(
 \|\cD^{\rm mix}_{E,\rho}x\|_{\mathcal Y_\tau^p}
 +\|x\|_{L^p}\bigr).
\end{equation}
The coefficient bounds needed to absorb the frozen errors follow from
Theorem~\ref{thm:inverse-adiabatic}.  The cross terms are triangular:
holomorphic variation enters the HYM and Coulomb rows, whereas unitary
gauge invariance and the K\"ahler identities prevent a gauge-to-HYM or
metric-to-Coulomb term.  Solving the harmonic row first and then the two
rows in \eqref{eq:two-block-estimates} gives uniformly bounded Schur
coordinates.

It remains to remove the last term in
\eqref{eq:uniform-dfl-with-compact-term}.  Otherwise there are
$E_j\to E_\infty$, $\rho_j\to0$, and reduced fields $x_j$ of unit graph
norm with $\cD^{\rm mix}_{E_j,\rho_j}x_j\to0$.  In the Schur coordinates,
the harmonic component converges to its Dolbeault kernel projection.
Orthogonality to the smoothly varying kernel makes that projection zero.
The two estimates in \eqref{eq:two-block-estimates} then make the metric and
gauge components tend to zero.  This contradicts the unit graph norm.
Thus
\[
 \|x\|_{\mathcal X_{E,\tau}^p}
 \leq C\|\cD^{\rm mix}_{E,\rho}x\|_{\mathcal Y_\tau^p}
\]
on the reduced complement.  Fixed-scale surjectivity from
Proposition~\ref{prop:fredholm-dolbeault} and the open-mapping theorem give
the right inverse and \eqref{eq:uniform-inverse}.
\end{proof}

On the seam-admissible locus, $G_\RR$ normalizes every gauge-side
basepoint.  A sequence whose basepoint approaches the seam is governed by
the center--scale ratio in \eqref{eq:rank-drop-exact-normal-form}: finite
ratio returns to the full-rank mixed chart, negative infinite ratio returns
to the normalized gauge rank-drop chart, and positive infinite ratio is the
ordinary holomorphic-sphere stratum treated in
Proposition~\ref{prop:ordinary-bubbles}.

\begin{proposition}[uniform attachment decay]\label{lem:attachment-decay}
Let $C$ be as in Theorem~\ref{lem:uniform-inverse}.  In a common seam gauge
near infinity there are constants $C_k$ such that
\begin{equation}\label{eq:attachment-decay}
 \sum_{j=0}^k|\zeta|^j
 \bigl|\nabla^j\bigl((A_{E,\rho},u_E)
                   -(A_\infty,u_\infty)\bigr)\bigr|
 \leq C_k|\zeta|
\end{equation}
for $E\in C$ and $0<\rho\leq\rho_0$, where $\zeta=1/w$ is the removable
coordinate at the attaching point.
\end{proposition}

\begin{proof}
Theorem~\ref{thm:inverse-adiabatic} gives uniform smooth bounds up to the
boundary point $\zeta=0$.  In temporal seam gauge, the value there is the
Narasimhan--Seshadri representative of $E_\infty$.  The positive
classifying maps form a compact smooth family with the same value.
Taylor's theorem gives \eqref{eq:attachment-decay}; the powers of
$|\zeta|$ record the scale-invariant derivative norm.  For
$\rho$ bounded away from zero the same estimate follows from ordinary
compact parameter elliptic theory.  Proposition~\ref{prop:primitive-no-neck}
identifies these representatives with the clean cut-circle limits of an
arbitrary degenerating sequence.
\end{proof}

\begin{theorem}[primitive master compactness]
\label{thm:primitive-master-compactness}
Let $(A_\nu,u_\nu)$ be mixed solutions with one matching concentration
point of primitive charge and no ordinary instanton charge.  Normalize the
center by the energy barycenter, the scale by a fixed energy radius, and
the remaining seam-preserving affine freedom.  After passage to a
subsequence, precisely one of the following occurs.
\begin{enumerate}
\item The coefficient pairs converge at finite seam ratio to a full-rank
$E\in\frakM_g^{\rm sa}$.  On every fixed rescaled core, the solutions
converge in the graph norm of Theorem~\ref{lem:uniform-inverse} to
$(A_{E,\rho_\nu},u_E)$.
\item After negative affine recentering, they converge to a rank-drop
$E\in\frakM_g^{\rm sa}$ whose basepoint lies in $H_-$.  The same strong
convergence holds to its Dirichlet reconstruction.
\item The seam ratio tends to positive infinity.  The seam-centered
component converges strongly to the constant mixed solution, while
recentering at the escaping basepoint gives a unique simple minimal sphere.
The convergence there is smooth away from the node.
\end{enumerate}
In every case the annulus from the primitive core to the residual solution
has vanishing limiting energy.  There is no further analytic scale or
positive-energy component.
\end{theorem}

\begin{proof}
Apply finite-scale holomorphic doubling before taking a limit.  On the
rescaled discs this gives bundles
\[
                     E_\nu\longrightarrow D_{R_\nu}\times\Sigma_g,
                     \qquad R_\nu\longrightarrow\infty,
\]
framed on regular outer circles chosen from smooth convergence to the
removable residual solution.  The classifying loop there lies in a
contractible fine-moduli chart and converges to the attaching value.
Cap it continuously in that chart.  This cap is used only for relative
degree; it is not asserted to be holomorphic.  It may be chosen with
symplectic area tending to zero.

Apply Lemma~\ref{lem:disc-reduction} on each expanding disc.  The
fixed-family identity \eqref{eq:finite-charge-splitting} gives
\[
                         d_\nu+\ell_\nu=1.
\]
The stable-reduction map is holomorphic on the disc, so its area is
$d_\nu/2+o(1)$, the error being the area of the shrinking cap.  Gromov
compactness followed by removal at infinity produces a stable map on
$\PP^1$, and the finitely many modification supports have a convergent
subsequence.  Positivity and the displayed identity leave either one
degree-one component or one length-one modification.  Mu\~noz's
classification identifies the result with one primitive coefficient pair.
The exact normal form \eqref{eq:rank-drop-exact-normal-form} then gives the
finite-ratio, negative-rank-drop, and positive-sphere alternatives.  It
also shows that a sphere and the limiting quasimap basepoint are two
compactifications of the same unit of charge, not two components to be
counted separately.

Consider the first two alternatives and fix a rescaled core $K_R$.  For
large $\nu$, stable reduction supplies complex gauges in which the
Dolbeault operators converge smoothly to that of $E$ on $K_R$.  The seam
Narasimhan--Seshadri metrics converge smoothly because all seam fibers are
stable and their trace-free zero-form Laplacians have a uniform spectral
gap.  Compare the original rescaled HYM metric with the reconstructed
metric in these gauges.  Their seam values agree in the limit.  On the
outer semicircle their difference tends to zero by residual convergence
and Proposition~\ref{lem:attachment-decay}.

Extend that outer boundary difference through a collar and subtract it.
If $e^{r_\nu}$ is the remaining metric ratio, then $r_\nu$ has zero
boundary value and satisfies an integrated linearized HYM equation
\begin{equation}\label{eq:compactness-metric-ratio}
 \left(\int_0^1
 D_H\Phi(H_{\rm rec}e^{t r_\nu},\rho_\nu^2)\,dt\right)r_\nu
 =f_\nu,
 \qquad \|f_\nu\|_{\mathcal Y_{\rho_\nu^2}^p(K_R)}\longrightarrow0.
\end{equation}
The maximum principle first gives $r_\nu\to0$ in $C^0$.  The metric and
gauge block estimates in \eqref{eq:two-block-estimates} then give
\[
 \|r_\nu\|_{\mathcal M_{\rho_\nu^2}^p(K_R)}
 +\|\chi_\nu\|_{\mathcal G_{\rho_\nu^2}^p(K_R)}\longrightarrow0,
\]
where $\chi_\nu$ is the unitary difference in the Coulomb slice.  The
harmonic part is precisely the convergent Dolbeault coefficient.  The
collar decomposition \eqref{eq:dfl-collar-hodge} therefore proves strong
convergence in the full graph norm.  The argument is uniform through
gauge-side rank drop because the unstable fiber stays in the interior;
the seam spectral gap and base Dirichlet estimate do not degenerate.

In the third alternative, the exact coefficient normal form converges on
seam-centered compact sets to the constant stable bundle, so the preceding
argument gives the constant mixed limit there.  In the coordinate centered
at the escaping basepoint the positive maps converge to
\[
                            \eta\longmapsto[v\eta+w].
\]
This is a simple minimal sphere.  Its linearized operator is surjective by
the regularity argument in Section~\ref{sec:compact-stratification}; removal
and elliptic bootstrapping give smooth convergence away from the node.

Finally let $R\to\infty$ along the good circles in
Proposition~\ref{prop:primitive-no-neck}.  Its Chern--Simons identity makes
the intervening energy tend to zero, and small-energy regularity upgrades
this to first-order convergence on the selected annuli.  The primitive
charge has already been exhausted.  A further matching component or sphere
would cost another half-unit, while an instanton would cost one unit.
This proves exhaustion and the theorem.
\end{proof}

\subsection{The primitive gluing collar}\label{sec:primitive-gluing}
Let $x_0=(A_0,u_0)$ be a regular index-zero residual mixed solution and
let
\[
                         e_0:\RR\longrightarrow M_g
\]
be its evaluation along the matching line.  The primitive bubble stratum
over $x_0$ is
\begin{equation}\label{eq:bubble-stratum}
 \cB_g(x_0)=
 \left(
 \RR\,{}_{e_0}\!\times_{\ev_\infty}\frakM_g^{\rm sa}
 \right)/G_\RR.
\end{equation}
When the fiber product is regular, the dimension calculation
\eqref{eq:bubble-dim-ledger} gives
\begin{equation}\label{eq:bubble-stratum-dim}
                           \dim_\RR\cB_g(x_0)=3.
\end{equation}
For $g\geq3$ the generic fiber product is empty by
\eqref{eq:image-codim}; the construction below applies on every regular
nonempty stratum before that final generic choice.

\begin{theorem}[primitive mixed gluing]\label{thm:primitive-gluing}
Suppose $x_0$ and the fiber product in \eqref{eq:bubble-stratum} are
regular.  For every compact $C\subset\cB_g(x_0)$ there are $\rho_0>0$
and a smooth map
\begin{equation}\label{eq:gluing-map}
 \Glue:C\times(0,\rho_0)
       \longrightarrow\cM_\eta(\alpha,\beta)_4
\end{equation}
with the following properties.
\begin{enumerate}
\item As $\rho\to0$, it converges in the
Gromov--Uhlenbeck--stable-reduction topology to the corresponding
residual-plus-bubble configuration.
\item It is a local diffeomorphism onto a neighborhood of the bubbled face.
\item Every mixed solution sufficiently close to $C$ is represented
uniquely by \eqref{eq:gluing-map}, modulo gauge and $G_\RR$.
\end{enumerate}
Thus $\rho=0$ is a collar face of the compactified index-four moduli space.
\end{theorem}

\begin{proof}
We give the construction with the estimates needed later.  In the
removable coordinate $\zeta=1/w$, the bubble and residual solutions have
the same stable flat seam value at $\zeta=0$.  Choose
\begin{equation}\label{eq:gluing-radii}
                      R_\rho=\rho^{-1/2},
 \qquad r_\rho=\rho R_\rho=\rho^{1/2},
 \qquad T_\rho=\frac12|\log\rho|.
\end{equation}
Cut the bubble at $|w|=R_\rho$, cut the residual solution at physical
radius $r_\rho$, and identify the annuli by $z=\rho w$.  Use a common
cutoff which is affine in logarithmic radius on a subcylinder of length
comparable to $T_\rho$.  Its cylindrical derivative is at most
$C/T_\rho$.  Use the same cutoff on both sides of the seam and the same
flat representative.  The matching condition is preserved exactly.  The
secondary perturbations vanish on this neighborhood and introduce no
rescaled error.

The residual Taylor error is $O(r_\rho)$, the bubble-tail error is
$O(R_\rho^{-1})$ by Proposition~\ref{lem:attachment-decay}, and the cutoff
commutator has $L^p$ size $O(T_\rho^{-\kappa})$, where
$\kappa=1-1/p>0$: its pointwise size is $O(T_\rho^{-1})$ and its transition
cylinder has length $O(T_\rho)$.  In the target norm of
Theorem~\ref{lem:uniform-inverse},
\begin{equation}\label{eq:pregluing-error}
 \|\cF(x_{\rm pre})\|_{\mathcal Y_{\rho^2}^p}
 \leq C\bigl(\rho^{1/2}+T_\rho^{-\kappa}\bigr)
\end{equation}
uniformly over $C$.

Let $Q_0$ be a right inverse for the regular index-zero residual operator
and let $Q_{E,\rho}$ be the reduced inverse from
Theorem~\ref{lem:uniform-inverse}.  If $K_0$ and $K_E$ are the two kernels,
regularity of the evaluation fiber product is surjectivity of
\[
 \mathfrak e:\RR_s\oplus K_0\oplus K_E\longrightarrow TM_g,
 \qquad
 \mathfrak e(\dot s,\xi_0,\xi_E)
 =d\ev_0(\dot s,\xi_0)-d\ev_\infty(\xi_E).
\]
Choose a uniformly bounded right inverse for this finite-dimensional map
over $C$.  Cut a target field into residual, bubble, and neck pieces, apply
$Q_0$ and $Q_{E,\rho}$, correct their harmonic attaching mismatch with
that right inverse, and splice the source fields with the logarithmic
cutoff.  The harmonic neck modes are exactly those treated by
$\mathfrak e$; irreducibility removes trace-free exact zero modes.

The resulting approximate inverse satisfies
\[
 \|D_\rho Q_\rho^{\rm app}-\id\|
 \leq C\bigl(T_\rho^{-1}+\rho^{1/2}\bigr).
\]
The cutoff term is small because its logarithmic derivative tends to zero,
the coefficient term is controlled by
Proposition~\ref{lem:attachment-decay}, and the collapsing block is
controlled by Theorem~\ref{lem:uniform-inverse}.  A Neumann series gives a
right inverse $Q_\rho$ with
\begin{equation}\label{eq:glued-inverse}
                              \|Q_\rho\|\leq C.
\end{equation}

The nonlinear term obeys the standard quadratic estimate
\[
 \|\mathcal N(\xi)-\mathcal N(\xi')\|
 \leq C(\|\xi\|+\|\xi'\|)\|\xi-\xi'\|.
\]
The contraction equation
\[
 \xi=-Q_\rho\bigl(\cF(x_{\rm pre})+\mathcal N(\xi)\bigr)
\]
therefore has a unique solution satisfying
\[
            \|\xi\|\leq C(\rho^{1/2}+T_\rho^{-\kappa}).
\]
It varies smoothly in the bubble parameters and $\rho>0$ and defines
\eqref{eq:gluing-map}.

For the boundary smooth structure use
\[
 r=T_\rho^{-1}=\frac{2}{|\log\rho|},
 \qquad \rho=e^{-2/r}.
\]
In this profile every positive power of $\rho$ is flat at $r=0$, and the
cutoff is fixed after dividing the neck coordinate by $T_\rho$.  The
parameter-dependent contraction theorem makes the corrected family smooth
for $r\geq0$.  Increasing $r$ has the same inward orientation as increasing
$\rho$.

For uniqueness, fix center and scale by the energy barycenter and the
half-energy radius and impose Coulomb gauge.  The derivative of these two
conditions on translation and positive dilation is nonsingular.  Two glued
solutions with the
same limit differ in the reduced complement; \eqref{eq:glued-inverse} and
the quadratic estimate force the difference to vanish.

Conversely, Theorem~\ref{thm:primitive-master-compactness} assigns every
nearby solution a residual solution, primitive coefficient, center, and
scale, and makes it converge to the corresponding pregluing in the source
graph norm on both cores and the complete neck.  It therefore lies in the
same contraction ball as the correction constructed above.  Uniqueness of
the fixed point proves local surjectivity without assuming uniqueness of an
abstract stable-reduction limit.

Changing cutoffs, gauges, or exponential charts changes the pregluing by a
term tending to zero in the source norm.  The two corrections solve the
same equation in the same center--scale slice, so contraction uniqueness
identifies them.  Coefficient-chart overlaps are smooth by Dirichlet
reconstruction, including at gauge rank drop.  At positive infinite ratio,
the exact coefficient normal form is the ordinary nodal-sphere smoothing;
standard simple-sphere gluing gives the compatible raw compactification
chart.  This proves the overlap assertion and completes the collar theorem.
\end{proof}

A rank-drop basepoint on the gauge side produces no additional end:
$G_\RR$ normalizes it, so it belongs to the same Dirichlet family as the
full-rank bubbles.  A positive-side rank-drop limit is instead a constant
mixed component with an ordinary minimal sphere.  Its standard nodal sphere
collar is part of the raw compactification and is compatible with the
finite seam-ratio chart through
\eqref{eq:rank-drop-exact-normal-form}; Proposition~\ref{prop:ordinary-bubbles}
shows that the point incidence misses it.  Thus
Theorem~\ref{thm:primitive-gluing} is required only on the full-rank and
gauge rank-drop colors of the primitive matching face.

\section{Transversality, incidence, and the boundary formula}
\label{sec:boundary-formula}

\subsection{Compact stratification}\label{sec:compact-stratification}
Through Fredholm index four, the compactified mixed moduli space has the
following strata:
\begin{enumerate}
\item the unbroken, unbubbled mixed stratum;
\item ordinary gauge and symplectic Floer breakings;
\item interior sphere and outer-boundary disc strata, which have
codimension at least two in the point cut by
Proposition~\ref{prop:ordinary-bubbles};
\item the primitive matching face constructed in
Theorem~\ref{thm:primitive-gluing};
\item the full-rank and Segre rank strata inside that face.
\end{enumerate}
There are no further scale-tree strata.  Theorem~\ref{thm:local-charge} and
Corollary~\ref{cor:primitive-atom} give one charge-one atom;
Proposition~\ref{prop:framed-splitting} gives each individual doubled
family's algebraic type as a minimal line or one rank-drop modification;
Proposition~\ref{prop:primitive-no-neck}
excludes hidden analytic necks; and Theorem~\ref{thm:primitive-gluing}
constructs a collar of the resulting face.

The obstruction to extending DFL's perturbation construction verbatim
from index three to index four is now visible.  Their local
cokernel-separation argument is index independent, but their finite-cover
argument needs compactness of the nonregular locus.  At index four that
compactness becomes available only after the primitive face has been
added.

The ordinary bubbling strata cause no further transversality problem.
Index exhaustion makes every sphere or disc component minimal and hence
simple.  Away from the matching neighborhood, the usual universal
almost-complex-structure argument regularizes these components and their
nodal evaluation maps.  In the region where the complex structure is the
fixed structure $J_*$, a minimal sphere is one of the extension lines in
$M_g$.  For its bundle $E$ on $\PP^1\times\Sigma_g$, stability of the
fibers gives $R^0\pi_*\End_0E=0$ and identifies
$R^1\pi_*\End_0E$ with the pullback of $TM_g$.  The Leray sequence and
the vanishing in Proposition~\ref{prop:fredholm-dolbeault} therefore give
$H^1(\PP^1,f^*TM_g)=0$, which is surjectivity of the sphere operator.  A
minimal outer-boundary disc is regular for a
generic allowed domain-dependent almost complex structure.  The standard
simple-curve gluing construction then gives neighborhoods of these
codimension-two strata; for the sphere case see
\cite[Chapter~10]{McDuffSalamon04}, and for the monotone Lagrangian
boundary case see \cite{Oh93}.  These
ordinary neighborhoods are used below only to compactify the
transversality problem; Proposition~\ref{prop:ordinary-bubbles} shows that
the point cut misses them.

\subsection{Evaluation transversality}\label{sec:evaluation-transversality}
Let $\cP$ be the DFL space of standard secondary perturbations: holonomy
perturbations on the gauge side and domain-dependent almost-complex
variations on the symplectic side, supported away from the matching
neighborhood.

Let $\cS$ be a stratum in an algebraic Whitney stratification for which
$i_g$ and the restrictions of $\ev_\infty$ to the full-rank and Segre
strata are stratified submersions.  On the universal residual moduli space
consider
\begin{equation}\label{eq:universal-seam-evaluation}
 \Ev_\gamma:\cM^{\rm univ}_0(\alpha,\beta)\times\RR_s\longrightarrow M_g,
 \qquad
 \Ev_\gamma([A,u],\eta,s)=[A|_{(s,0)\times\Sigma_g}].
\end{equation}
The matching condition gives the same value from $u$.

\begin{proposition}[relative seam evaluation]\label{lem:seam-evaluation}
After an arbitrarily small finite-dimensional enlargement of $\cP$, the
map \eqref{eq:universal-seam-evaluation} is transverse to every stratum
$\cS$ occurring in the compact energy range.  The enlargement is
supported away from the matching neighborhood and can be made relative to
all previously regular moduli spaces.
\end{proposition}

\begin{proof}
Write $\cF_\eta$ for the mixed equation and fix a point at which
$\Ev_\gamma$ lies in $\cS$.  Transversality of
$(\cF,\Ev_\gamma)$ to the zero section times $\cS$ is equivalent to
surjectivity of
\begin{equation}\label{eq:augmented-evaluation-linearization}
 (\xi,\dot\eta,\dot s)\longmapsto
 \left(
 D_{(A,u)}\xi+P_{(A,u)}\dot\eta,
 \pi_{N\cS}\,d\Ev_\gamma(\xi,\dot s)
 \right).
\end{equation}
Suppose a pair $(\lambda,\nu)$ annihilates its image.  Here $\lambda$ is
an adjoint mixed field and $\nu\in N^*\cS$.  Variation in $\dot\eta$
shows that $\lambda$ annihilates every allowed holonomy and
almost-complex direction.  On the complement of the evaluation point it
satisfies DFL's homogeneous adjoint equation.  If its gauge component is
nonzero on the allowed support, the holonomy direction constructed in the
proof of \cite[Proposition~3.23]{DaeFukLip21a} pairs nontrivially with it.
If the gauge component vanishes but the symplectic component is nonzero,
DFL's variation of the almost complex structure at a point where $du$ is
nonzero does the same.  Their constant-map argument covers the remaining
case.  Thus $\lambda$ vanishes on the perturbation supports.  Unique
continuation on both pieces makes it vanish away from the evaluation
point.

The weak adjoint equation is now
\begin{equation}\label{eq:adjoint-evaluation-source}
                       D_{(A,u)}^*\lambda+\Ev_\gamma^*\nu=0.
\end{equation}
The adjoint field belongs to the ordinary dual Sobolev space, so vanishing
off one point means that $\lambda=0$ as an $L^q$ field.  Equation
\eqref{eq:adjoint-evaluation-source} reduces to $\Ev_\gamma^*\nu=0$.
Smooth configuration variations realize an arbitrary tangent vector in
$TM_g$ at the seam trace: prescribe the matching trace locally and extend
it to the two configuration spaces.  Hence $\nu=0$, proving surjectivity of
\eqref{eq:augmented-evaluation-linearization}.

The relevant universal zero set is compact after the ordinary and
primitive strata have been added.  Finitely many of the separating
directions above therefore suffice.  Sard--Smale applied to this
finite-dimensional family gives the assertion.  The usual descending
smallness choice makes the construction relative to earlier regular
spaces and preserves DFL's nonnegative-energy bound.
\end{proof}

Choose a Whitney stratification of the algebraic image of $i_g$ and apply
Proposition~\ref{lem:seam-evaluation} simultaneously to its finitely many strata
in the compact energy range.  For $g\geq3$, the residual path misses the
image by \eqref{eq:image-codim}.  For $g=2$, it avoids the branch locus and
meets the regular four-sheeted locus.  When $g=1$, the target is a point;
the excess Jacobian is regularized by the universal incidence section
rather than by evaluation.

\begin{proposition}[relative index-four regularity]
\label{prop:index-four-regularity}
The secondary perturbations can be chosen so that all mixed moduli spaces
needed by the point operation are regular through index four, including
their residual, broken, primitive, and rank strata.  The perturbations
remain identically zero in the fixed matching neighborhood.
\end{proposition}

\begin{proof}
First choose the DFL perturbations so that all lower-index residual and
broken strata are regular.  Proposition~\ref{lem:seam-evaluation} makes
the primitive evaluation fiber products regular.  Fixed-scale regularity
and Theorem~\ref{thm:primitive-gluing} then give parameterized collars of
the primitive faces.  Ordinary breaking collars are regular by DFL Floer
gluing, and the simple sphere and disc collars are regular by the preceding
ordinary-bubble argument.

These constructions also apply to a sufficiently small family of
secondary perturbations: the right inverses and gluing maps vary
continuously, so regularity of the compact union of boundary strata is an
open condition.  Remove smaller collar neighborhoods from the
parameterized compactified index-four space.  Properness of DFL compactness
together with Theorem~\ref{thm:primitive-gluing} makes the remaining
universal core compact.  At each nonregular solution in this core, DFL's
local cokernel-separation argument supplies finitely many perturbation
directions that kill its cokernel.  A finite subcover and Sard's theorem,
exactly as in the proof of \cite[Proposition~3.23]{DaeFukLip21a}, give one
arbitrarily small parameter for which the core is regular.  Smallness
preserves the already regular collar strata.  All perturbation directions
are supported away from the seam, so the local charge, stable-reduction,
Dirichlet, and primitive gluing models are unchanged.
\end{proof}

\subsection{Extension of the point incidence}\label{sec:incidence-extension}
Blow up the pair $(z-z_0,\rho)$ near a primitive matching point.  The front
face has affine coordinate $w=(z-z_0)/\rho$.  In the common flat seam gauge
on a good-circle annulus, let $g_0$ and $g_E$ be the framed transition
functions of the residual and reconstructed bubble bundles.  For
$\rho>0$, splice $g_0^{-1}g_E$ with the logarithmic cutoff used in the
linear gluing construction.  At $\rho=0$, retain the residual cocycle on
the principal component and the bubble cocycle on the front face, identified
at $w=\infty$.  No-neck convergence makes the cocycles converge on the
overlap.  This constructs the universal projective bundle on the gluing
blowup.

The fine-moduli ambiguity tensors a rank-two universal bundle by a line
from the parameter space; its adjoint is unchanged.  We therefore obtain a
canonical adjoint bundle over the compactified chart, up to an isomorphism
fixed at the interpolation ends.  It is smooth on every stratum and through
gauge rank drop by the uniform reconstruction theorem.  On the positive
infinite-ratio chart it is the standard universal bundle over the ordinary
stable-map sphere collar.  Its front-face clutching class is the class used
in Theorem~\ref{thm:local-charge}, so there is no additional neck class.

Two sections of the extended bundle define a stratified codimension-four
linear-dependence cycle on
\[
               \overline{\cM}_\eta(\alpha,\beta)_4\times\RR_t.
\]

\begin{proposition}[relative point-incidence transversality]
\label{prop:incidence-transversality}
The two interpolating sections can be chosen, relative to the fixed
transverse representatives at $t=\pm\infty$, so that their rank-one
degeneracy locus is transverse to every stratum of the index-three and
compactified index-four universal evaluation maps.  The rank-zero locus is
empty.  On the genus-one primitive stratum the induced normal section of
the excess determinant line is transverse to zero.
\end{proposition}

\begin{proof}
At a point of the rank-one stratum, independent pointwise variations of
the two sections surject onto the normal space of the determinantal variety
in $\Hom(\CC^2,\CC^3)$.  The same statement, with a larger normal space,
holds at the rank-zero stratum.  Section variations supported in an
arbitrarily small neighborhood of a given evaluation point realize these
pointwise variations.  Relative Thom transversality, applied successively
to the finite Whitney stratification from
Proposition~\ref{prop:index-four-regularity}, therefore gives simultaneous
transversality while leaving collars of $t=\pm\infty$ fixed.  The
rank-zero stratum has real codimension twelve, greater than the dimension
of every evaluation stratum here, and is consequently empty.  On the
genus-one excess family, prescribe the first complex normal derivative as
a section of the rank-two coefficient bundle.  Its real rank is four while
the base has real dimension two, so relative transversality makes it
nowhere zero.  Project the second normal derivative to the resulting
complex cokernel line.  A further relative variation makes that section
transverse to zero.  These are the two derivatives used in
Proposition~\ref{prop:algebraic-local-determinant}.
\end{proof}

Near a primitive face, \eqref{eq:bubble-stratum-dim} and the interpolation
parameter give dimension four.  Proposition~\ref{prop:incidence-transversality}
cuts the face in an oriented zero-manifold and its gluing collar in a
one-manifold.
Together with Proposition~\ref{prop:ordinary-bubbles} and
Proposition~\ref{prop:index-four-regularity}, this gives the required
compactness statement.

\begin{proposition}[compact point-cut moduli spaces]
\label{prop:point-manifolds}
For compatible generic choices, $\cL_0^q(\alpha,\beta)$ is an oriented
zero-manifold and $\overline{\cL}_1^q(\alpha,\beta)$ is a compact oriented
one-manifold.  Its finite-$t$ boundary consists only of ordinary Floer
breakings and primitive matching configurations.
\end{proposition}

\begin{proof}
Propositions~\ref{prop:index-four-regularity} and
\ref{prop:incidence-transversality} give the asserted dimensions and
orientations on the top stratum and on each compactification stratum.  The
primitive gluing collar makes the primitive cut compact, while
Proposition~\ref{prop:ordinary-bubbles} excludes the ordinary sphere and
disc faces.  DFL compactness, index-four exhaustion, and primitive stable
reduction leave only the boundary types listed in the statement.  The
index-three cut is bubble-free by Proposition~\ref{prop:index-exhaustion};
its zero-dimensional compactification therefore has no additional ends.
\end{proof}

\subsection{Determinant locality}\label{sec:determinant-locality}
Orient $G_\RR$ by translation followed by positive dilation, orient the
matching-center and interpolation lines $\RR_s,\RR_t$ by increasing
coordinate, and write $\mathfrak g_\RR$ for the oriented Lie algebra of
$G_\RR$.  Let $\mathcal N_{p_1}$ be the complex rank-two normal bundle of
the rank-one determinantal stratum in
$\Hom(\CC^2,\CC^3)$, with DFL's positive $p_1$ coorientation.

\begin{proposition}[determinant excision]\label{prop:determinant-excision}
At a regular primitive configuration $(x_0,E)$, the determinant line of
the glued index-four operator has a natural excision isomorphism
\begin{equation}\label{eq:det-excision}
 \det\cD_{\rm gl}
 \cong
 \det\cD_{x_0}\otimes
 \det\cD_{E,\rho}^{\rm mix}\otimes
 \det\RR_s\otimes
 (\det TM_g)^*\otimes
 (\det\mathfrak g_\RR)^*\otimes
 \RR_\rho.
\end{equation}
Here $\RR_\rho$ is the gluing-normal line, oriented by increasing
$\rho$.  The isomorphism is continuous across the full-rank and rank-drop
strata.  On the boundary point cut, its finite-dimensional local factor is
\begin{equation}\label{eq:local-orientation-line}
 \begin{split}
 \mathfrak o_g(E)={}&
 \det\cD_{E,\rho}^{\rm mix}\otimes
 \det\RR_s\otimes(\det TM_g)^*\otimes
 (\det\mathfrak g_\RR)^*\\
 &\otimes\det\RR_t\otimes(\det\mathcal N_{p_1})^*.
 \end{split}
\end{equation}
Thus the orientation line of a primitive point-cut boundary point is
$\det\cD_{x_0}\otimes\mathfrak o_g(E)$.
\end{proposition}

\begin{proof}
The patched inverse in Theorem~\ref{thm:primitive-gluing} deforms the
linearized glued operator, through Fredholm operators, to the direct sum of
the residual and fixed-domain bubble operators together with the
finite-dimensional matching map.  At the level of kernels that map is
\begin{equation}\label{eq:matching-linear-map}
 \RR_s\oplus H^1(\PP^1\times\Sigma_g,\End_0E)
 \longrightarrow TM_g,
 \qquad
 (\dot s,\dot E)\longmapsto
 de_0(\dot s)-d\ev_\infty(\dot E).
\end{equation}
It is onto precisely when the evaluation fiber product
\eqref{eq:bubble-stratum} is regular.  Proposition~\ref{lem:seam-evaluation}
provides this condition after the generic relative choice of perturbations.
Its kernel contains
the infinitesimal $G_\RR$ orbit, and the tangent of the primitive face is
the quotient by that orbit.  The determinant functor applied to the two
resulting short exact sequences gives all factors in
\eqref{eq:det-excision}.  The gluing parameter supplies $\RR_\rho$.

Proposition~\ref{prop:fredholm-dolbeault} identifies the bubble determinant
with the complex determinant of the Dolbeault complex and shows that its
obstruction space vanishes.  The matching target $TM_g$ has its complex
orientation.  The constructions depend continuously on $E$, including on
the Segre locus, because the patched inverses and the Dolbeault family do.
Finally, adding $\RR_t$ and applying the determinant functor to the
surjective normal derivative of the point incidence contracts
$\det\mathcal N_{p_1}$ and gives \eqref{eq:local-orientation-line}.
\end{proof}

All factors in $\mathfrak o_g(E)$ are supported in the standard
surface-seam neighborhood.  They are independent of the two
three-manifolds and of the residual operator; the latter contributes only
$\det\cD_{x_0}$, with the orientation used in the comparison count $N$.
Consequently the primitive contribution is multiplication of the
comparison count by a universal integer
\begin{equation}\label{eq:local-coefficient}
                              c_g^{\rm loc}\in\ZZ
\end{equation}
depending only on the genus, the integral $p_1$ convention, and DFL's
seam orientation.  Summing over index-zero residual solutions gives
\begin{equation}\label{eq:B-local}
                              B_q=c_g^{\rm loc}N.
\end{equation}

\begin{proposition}[algebraic local determinant]\label{prop:algebraic-local-determinant}
With the complex orientations in
Proposition~\ref{prop:fredholm-dolbeault}, the absolute local count is four
in genus two.  In genus one the excess local determinant line over
$J(\Sigma_1)$ is canonically
\begin{equation}\label{eq:torus-det-identification}
                              \det\cE_\zeta.
\end{equation}
Hence its absolute local count is eight.
\end{proposition}

\begin{proof}
The Dolbeault reduction makes the finite-dimensional normal problem
complex linear.  In genus two, evaluation at infinity is the holomorphic
map $i_2:N_2\to M_2$.  Over a regular value, the remaining normalized
affine, center, and interpolation variables form the standard complex
rank-two normal slice to the two-section determinantal cycle.  To see its
relative degree, fix an extension basis $(v,w)$ with $v$ representing the
attaching line.  A nearby pencil is
\[
                  [v z_0+(a v+c w)z_1].
\]
The complex number $c$ is the Segre smoothing coordinate.  The matching
center and normalized interpolation displacement form a second complex
coordinate $\zeta$; the real affine quotient removes the complementary
translation and positive-dilation directions.  Before applying relative
transversality on the complement, prescribe the two-section pair on a
tubular neighborhood of this slice to be the Thom model whose
incidence-normal map is
\[
                         (c,\zeta)\longmapsto(c,\zeta)
\]
on every small normal disk.  It is nonzero on the disk boundary and has one
transverse zero of complex multiplicity one.  The pair extends over the
complement by Proposition~\ref{prop:incidence-transversality}.  Any two such
Thom choices are homotopic through nonzero boundary maps and have the same
degree.  The ordinary positive-sphere chart has
negative expected incidence dimension, while the map extends transversely
across the real-codimension-two gauge rank-drop stratum, so no degree escapes
through the compactification.  Thus every point of $i_2^{-1}(y)$ has
complex multiplicity one.  Proposition
\ref{prop:degree-four} gives four such points counted with multiplicity.

For $g=1$, the character variety is a point and every projective coefficient
pair has rank one.  After the basepoint and the real affine orbit are normalized,
the remaining excess parameter is $L\in J(\Sigma_1)$.  Applying the
determinant functor to the Dolbeault normal complex and to the rank-one
incidence normal sequence cancels the fixed domain, affine, and Dirichlet
metric factors.  More explicitly, the normal incidence complex over $L$
is
\begin{equation}\label{eq:torus-coefficient-complex}
 \CC_{\rm norm}\longrightarrow
 H^0(\PP^1,\cO(1))\otimes(\cE_1)_L.
\end{equation}
The source is the trivial complex line obtained from the center and
interpolation normal after the affine quotient; the target records the two
extension coefficients.  Use the relative freedom in the first normal
derivative of the incidence pair to make this arrow a nowhere-zero smooth
bundle map.  Its cokernel is then a complex line bundle $Q$ fitting into
\[
 0\longrightarrow\CC_{\rm norm}\longrightarrow\cE_\zeta
   \longrightarrow Q\longrightarrow0.
\]
The remaining normal derivative of the incidence pair projects to a
section $\sigma_Q$ of $Q$.  A primitive configuration lies in the point cut
exactly when $\sigma_Q$ vanishes, and
Proposition~\ref{prop:incidence-transversality} makes this section
transverse to zero.  Since the source line is canonically trivial,
\[
                  Q\cong\det\cE_\zeta,
\]
which is \eqref{eq:torus-det-identification}.  This calculation is
unaffected by the metric block because that block is an invertible family
with its canonical determinant trivialization.  Equation
\eqref{eq:torus-degree-eight} gives
\begin{equation*}
 \left\langle c_1(\det\cE_\zeta),[J(\Sigma_1)]\right\rangle=8.
\end{equation*}
The complex zero-set orientation gives the asserted absolute counts.
\end{proof}

\subsection{Evaluation of the local coefficient}\label{sec:coefficients}
For $g\geq3$, Proposition~\ref{lem:seam-evaluation} and
\eqref{eq:image-codim} make the attachment fiber product empty.  Hence
\begin{equation}\label{eq:higher-coefficient}
                           c_g^{\rm loc}=0,
                    \qquad g\geq3.
\end{equation}

Proposition~\ref{prop:algebraic-local-determinant} gives absolute local
counts four and eight in genera two and one, respectively.

It remains to compare these complex orientations with DFL's coherent mixed
orientation.  The excision isomorphism \eqref{eq:det-excision} proves that
this comparison is the same universal sign on every target.  Evaluate it
on the cylinder splitting of $\Sigma_g\times S^1$.  DFL orient their
index-$8k+4$ gauge path by gluing the positive generator sphere to a
constant strip and require the symplectic orientations to agree with that
sphere-gluing map \cite[Remark~4.76]{DaeFukLip21a}.  Independently,
Proposition~\ref{prop:cylinder-identities} gives $+8N$ and $+4N$ from the
Floer and quantum ring relations.  Since $N$ is an isomorphism, the local
complex orientations agree with the DFL boundary orientation.  Therefore
\begin{equation}\label{eq:integral-coefficients}
 c_g^{\rm loc}=
 \begin{cases}
 8,&g=1,\\
 4,&g=2,\\
 0,&g\geq3.
 \end{cases}
\end{equation}
This calibration is not circular: it uses surface compatibility, not
point compatibility.

\subsection{The oriented boundary count}\label{sec:oriented-boundary}

\begin{theorem}[point-class comparison]\label{thm:main}
Let $(Y_\#,E_\#)$ have an admissible splitting
$Y_\#=Y\cup_\Sigma(-Y')$ in the embedded-Lagrangian DFL setting.  Assume
the flat connections are irreducible and use compatible standard
perturbations that vanish near the matching line.  Let
$q\in\Sigma_j$ and write $g_j$ for the genus of that component.  Over
$\QQ$, the comparison chain map and the positive quarter-Pontryagin point
operations satisfy
\begin{equation}\label{eq:main}
 N m_q^G-m_q^S N
 =d_SK_q+K_qd_G+\kappa_{g_j}N,
 \qquad
 \kappa_g=
 \begin{cases}
 2,&g=1,\\
 1,&g=2,\\
 0,&g\geq3.
 \end{cases}
\end{equation}
\end{theorem}

\begin{proof}
The zero-dimensional index-three point cut defines $K_q$ by
\eqref{eq:K-integral}; it is bubble-free by
Proposition~\ref{prop:index-exhaustion}.  Proposition~\ref{prop:point-manifolds}
gives a compact oriented one-manifold at index four.

Its ordinary broken boundary points contribute
$d_SK_q+K_qd_G$.  Escape of the interpolation parameter to the gauge and
symplectic ends contributes $Nm_q^G-m_q^SN$ with DFL's end convention;
this is the same orientation calculation as
\cite[Proposition~7.12]{DaeFukLip21a}.  By
Proposition~\ref{prop:ordinary-bubbles}, no finite-$t$ interior sphere or
outer disc contributes.  Corollary~\ref{cor:primitive-atom},
Proposition~\ref{prop:framed-splitting}, and
Theorem~\ref{thm:primitive-gluing} show that the only remaining boundary
is the primitive matching face.  Determinant locality and
\eqref{eq:integral-coefficients} make its integral contribution
$c_{g_j}^{\rm loc}N$.

The signed boundary count of a compact oriented one-manifold is zero.
Using the integral identity \eqref{eq:boundary-with-B} and dividing by four
gives \eqref{eq:main}.
\end{proof}

The same construction with one additional continuation parameter proves
independence of the interpolation sections and secondary perturbations.
The parameterized compactification has the same primitive faces; the
local coefficient is unchanged by determinant locality, and its remaining
boundary gives the standard chain homotopy between choices.

\section{Integral normalization and consequences}
\label{sec:integral}

\begin{theorem}[integral point cycle]\label{thm:integral}
Before division by four, the operations represented by the positive
integral degeneracy cycle satisfy
\begin{equation}\label{eq:integral-main}
 N b_q^G-b_q^S N
 =d_SK_q^{p_1}+K_q^{p_1}d_G+c_{g_j}N,
 \qquad
 c_g=
 \begin{cases}
 8,&g=1,\\
 4,&g=2,\\
 0,&g\geq3.
 \end{cases}
\end{equation}
\end{theorem}

\begin{proof}
This is the oriented boundary count in the proof of
Theorem~\ref{thm:main} before the definitions
$m_q^G=\frac14b_q^G$, $m_q^S=\frac14b_q^S$, and
$K_q=\frac14K_q^{p_1}$ are applied.
\end{proof}

\begin{corollary}\label{cor:corrected-operator}
Define
\[
                  \widehat m_q^S=m_q^S+\kappa_{g_j}\id.
\]
Then $N$ intertwines $m_q^G$ and $\widehat m_q^S$ up to chain homotopy:
\begin{equation}\label{eq:corrected-intertwining}
             Nm_q^G-\widehat m_q^S N=d_SK_q+K_qd_G.
\end{equation}
In particular, the induced Atiyah--Floer isomorphism intertwines these
corrected point operations on homology.
\end{corollary}

\begin{proof}
Move $\kappa_{g_j}N$ to the left-hand side of
\eqref{eq:main} and use the definition of $\widehat m_q^S$.
\end{proof}

\subsection{The full corrected operation algebra}
\label{subsec:full-operation-algebra}

Theorem~\ref{thm:main} completes the list of generator comparisons.  We record
the resulting algebra statement, first in the integral normalization used by
the degeneracy-cycle constructions and then in the quarter-Pontryagin
normalization of the introduction.

Give $H_k(\Sigma;\ZZ)$ degree $4-k$.  For a commutative ring $R$, let
$V_R=H_*(\Sigma;\ZZ)\otimes_{\ZZ}R$ and define
\begin{equation}\label{eq:integral-observable-algebra}
 \mathbb A_{p_1}(\Sigma;R)
 =T_R(V_R)\big/
   \left\langle x\otimes y-(-1)^{|x||y|}y\otimes x\right\rangle
 =:\operatorname{Sym}^{\rm gr}_R(V_R)
\end{equation}
for homogeneous $x,y$.  This presentation fixes the meaning of
``graded-commutative'' over rings with two-torsion.  In particular, it
imposes $2x^2=0$, but not $x^2=0$, for an odd generator.  If $2$ is
invertible in $R$, then
\begin{equation}\label{eq:usual-observable-algebra}
 \mathbb A_{p_1}(\Sigma;R)
 \cong
 \operatorname{Sym}_R(H_0(\Sigma;R)\oplus H_2(\Sigma;R))
 \otimes_R\Lambda_R H_1(\Sigma;R).
\end{equation}

\begin{lemma}[integral operation relations]
\label{lem:integral-operation-relations}
On gauge and symplectic Floer homology with coefficients in $R$, the
integral universal-class operations satisfy
\[
 b_xb_y=(-1)^{|x||y|}b_yb_x
\]
for homogeneous $x,y\in V_R$.  Hence the homology actions factor through
\eqref{eq:integral-observable-algebra}.  For an odd loop class this asserts
$2b_x^2=0$; no vanishing of $b_x^2$ in characteristic two is used.
\end{lemma}

\begin{proof}
Work first over $\ZZ$.  The operation $b_x$ is the cap operation associated
with the slant class $p_1(\mathbb V)/x$ on the relevant configuration space;
DFL construct its gauge and symplectic representatives from the same
oriented integral degeneracy locus in Section~7.2 of
\cite{DaeFukLip21a}.  For two classes, use trajectories with two ordered
marked cutdowns.  Stretching the interval between the marked points gives
the composite $b_xb_y$.  Bringing the points together gives the cutdown by
the cup product
\[
 (p_1(\mathbb V)/x)\smile(p_1(\mathbb V)/y).
\]
The compactified one-parameter family has only its two endpoint cutdowns
and Floer breakings; the latter are the differential terms in the resulting
chain homotopy.  This is the usual cap-module proof on the gauge side and
the quantum-module proof on the symplectic side, with sphere terms included
in the quantum product.  Interchanging the two factors in the coincident
cutdown changes its orientation by $(-1)^{|x||y|}$.  Reversing the
deformation therefore proves the displayed relation on homology.

All parameter spaces, degeneracy cycles, and determinant orientations in
this argument are integral.  Tensoring the resulting chain homotopies with
$R$ proves the assertion over an arbitrary commutative ring.  Taking $x=y$
odd gives exactly $2b_x^2=0$.
\end{proof}

For $\sigma\in H_k(\Sigma;\ZZ)$, let $b_\sigma^G$ and $b_\sigma^S$
denote the gauge and symplectic operations obtained from the positive
integral degeneracy cycle $p_1/[\sigma]$.  The constructions are extended
$R$-linearly after changing coefficients.  Surface and loop operations are
left unchanged.  For $q\in\Sigma_j$, set
\begin{equation}\label{eq:integral-corrected-point}
 \widehat b_q^S=b_q^S+c_{g_j}\id,
 \qquad
 c_g=
 \begin{cases}
  8,&g=1,\\
  4,&g=2,\\
  0,&g\geq3.
 \end{cases}
\end{equation}
The resulting homology actions are denoted by $\rho_{G,p_1}^R$ and
$\widehat\rho_{S,p_1}^R$, and the scalar extension of $N$ is denoted by
$N_R$.

\begin{corollary}[full corrected operation algebra]
\label{thm:full-operation-algebra}
Under the hypotheses of Theorem~\ref{thm:main}, for every commutative ring
$R$ the base-changed comparison isomorphism is an isomorphism of
$\mathbb A_{p_1}(\Sigma;R)$-modules:
\begin{equation}\label{eq:full-integral-action}
 N_R\rho_{G,p_1}^R(a)
 =\widehat\rho_{S,p_1}^R(a)N_R,
 \qquad a\in\mathbb A_{p_1}(\Sigma;R).
\end{equation}
For each chosen word representing $a$, the equality is induced by an
explicit chain homotopy.  No strict chain-level graded-commutativity or
homotopy-coherence assertion is included.
\end{corollary}

\begin{proof}
Work first over $\ZZ$.  For a homogeneous generator $x$, write $A_x$ for
the gauge operation and $B_x$ for the corrected symplectic operation.  The
detailed DFL surface and loop cutdowns use the integral degeneracy locus
representing $p_1$.  Their comparison theorem therefore supplies maps $K_x$
for $x\in H_1(\Sigma)\oplus H_2(\Sigma)$ such that
\begin{equation}\label{eq:generator-graded-homotopy}
 NA_x-B_xN=d_SK_x+(-1)^{|x|}K_xd_G.
\end{equation}
For $x\in H_0(\Sigma)$, Theorem~\ref{thm:integral} gives the same identity
after the correction in \eqref{eq:integral-corrected-point} is moved to the
symplectic side.

Suppose \eqref{eq:generator-graded-homotopy} has been extended to homogeneous
words $u$ and $v$.  Define
\begin{equation}\label{eq:product-homotopy}
                 K_{uv}=K_uA_v+(-1)^{|u|}B_uK_v.
\end{equation}
Then
\begin{align*}
 NA_uA_v-B_uB_vN
 &=(NA_u-B_uN)A_v+B_u(NA_v-B_vN)\\
 &=d_SK_{uv}+(-1)^{|u|+|v|}K_{uv}d_G.
\end{align*}
The second equality is the graded Leibniz rule, using that $A_v$ and $B_u$
are chain maps of degrees $|v|$ and $|u|$.  Induction proves the identity for
every word, and linearity treats every polynomial expression.  Passing to
homology proves \eqref{eq:full-integral-action} for every ordered word.
Lemma~\ref{lem:integral-operation-relations} shows that the resulting
homology operators factor through \eqref{eq:integral-observable-algebra};
the construction of $K_a$ may depend on the chosen word representing $a$.

DFL prove that $N$ is an isomorphism of chain complexes over $\ZZ$.
Tensoring $N$, its inverse, the operation maps, and the displayed homotopies
with an arbitrary commutative ring $R$ preserves all identities.  No
flatness hypothesis on $R$ is required.
\end{proof}

\begin{corollary}[quarter-Pontryagin normalization]
\label{cor:full-quarter-operation-algebra}
Suppose $4$ is invertible in $R$ and put $m_\sigma=\frac14b_\sigma$.  Then
Corollary~\ref{thm:full-operation-algebra} holds for the normalized observable
action, with the corrected point class
\begin{equation}\label{eq:quarter-corrected-point}
 \widehat m_q^S=m_q^S+\kappa_{g_j}\id,
 \qquad
 (\kappa_1,\kappa_2,\kappa_{g\geq3})=(2,1,0).
\end{equation}
In particular, over $\QQ$ the DFL comparison intertwines the full corrected
Donaldson observable algebra
\begin{equation*}
 \operatorname{Sym}\bigl(H_0(\Sigma;\QQ)\oplus H_2(\Sigma;\QQ)\bigr)
 \otimes\Lambda H_1(\Sigma;\QQ).
\end{equation*}
\end{corollary}

\begin{proof}
Divide each integral generator identity and its homotopy by four.  The point
coefficients in \eqref{eq:integral-corrected-point} become respectively
$2$, $1$, and $0$.
\end{proof}

\begin{corollary}[operation images and kernels]
\label{cor:operation-images-kernels}
For every coefficient ring covered by
Corollary~\ref{thm:full-operation-algebra}, conjugation by $N_R$ identifies the
two image operation algebras.  Moreover,
\begin{equation}\label{eq:operation-kernels}
 \ker\rho_{G,p_1}^R=\ker\widehat\rho_{S,p_1}^R,
\end{equation}
so the identity on $\mathbb A_{p_1}(\Sigma;R)$ induces an isomorphism of the
two operation quotients.
\end{corollary}

\begin{proof}
Equation~\eqref{eq:full-integral-action} and invertibility of $N_R$ give
\eqref{eq:operation-kernels}; the same conjugacy identifies the images.
\end{proof}

\begin{remark}\label{rem:operation-algebra-boundary}
Corollary~\ref{thm:full-operation-algebra} concerns the action of the full
observable algebra.  It does not say that the mixed comparison map preserves
a gauge-theoretic pair-of-pants product, and the wordwise homotopies
\eqref{eq:product-homotopy} do not constitute an $A_\infty$ module morphism.
An integral operation representing $\frac14p_1/[\sigma]$ still requires the
divisibility data discussed in Remark~\ref{rem:coefficients}; the
all-coefficient statement instead uses the integral $p_1$ cycle.
\end{remark}

\begin{remark}\label{rem:uncorrected-false}
The correction in genera one and two cannot be removed by changing the
compactification.  Proposition~\ref{prop:cylinder-identities} detects it
algebraically on $\Sigma_g\times S^1$.  It is therefore part of the point
operation itself under the DFL comparison, rather than an artifact of the
proof.
\end{remark}

\begin{remark}\label{rem:coefficients}
The quarter-Pontryagin statement is asserted over $\QQ$.  An integral
operation representing $\frac14p_1/[q]$ requires a divisibility choice for
the universal bundle that is not used here.  The integral theorem is
Theorem~\ref{thm:integral}, formulated directly with the degeneracy cycle
$p_1/[q]$.
\end{remark}

\section*{Availability}
The manuscript source is available at
\begin{center}
\url{https://github.com/bwuebben/donaldson-atiyah-floer}.
\end{center}

\begin{thebibliography}{DFL21b}

\bibitem[AB83]{AtiyahBott83} M.~F. Atiyah and R.~Bott,
\emph{The Yang--Mills equations over Riemann surfaces}, Philos. Trans.
Roy. Soc. London Ser. A \textbf{308} (1983), no.~1505, 523--615.

\bibitem[BB16]{BaerBallmann16} C.~B\"ar and W.~Ballmann,
\emph{Guide to boundary value problems for Dirac-type operators},
Arbeitstagung Bonn 2013, Progr. Math., vol.~319, Birkh\"auser/Springer,
Cham, 2016, 43--80.

\bibitem[Ati88]{Atiyah88} M.~F. Atiyah, \emph{New invariants of
$3$- and $4$-dimensional manifolds}, The mathematical heritage of Hermann
Weyl (Durham, NC, 1987), Proc. Sympos. Pure Math. \textbf{48}, Amer. Math.
Soc., Providence, RI, 1988, 285--299.

\bibitem[BW18]{BottmanWehrheim18} N.~Bottman and K.~Wehrheim,
\emph{Gromov compactness for squiggly strip shrinking in pseudoholomorphic
quilts}, Selecta Math. (N.S.) \textbf{24} (2018), no.~4, 3381--3443.

\bibitem[DFL26]{DaeFukLip21a} A.~Daemi, K.~Fukaya, and M.~Lipyanskiy,
\emph{Lagrangians, $SO(3)$-instantons and the Atiyah--Floer conjecture},
accepted for publication, 2026; \texttt{arXiv:2109.07032}.

\bibitem[DFL24]{DaeFukLip21b} A.~Daemi, K.~Fukaya, and M.~Lipyanskiy,
\emph{Lagrangians, $SO(3)$-instantons and mixed equation},
Geom. Funct. Anal. \textbf{34} (2024), no.~3, 659--732.

\bibitem[Don92]{Donaldson92} S.~K. Donaldson, \emph{Boundary value
problems for Yang--Mills fields}, J. Geom. Phys. \textbf{8} (1992),
no.~1--4, 89--122.

\bibitem[DK90]{DonaldsonKronheimer90} S.~K. Donaldson and P.~B. Kronheimer,
\emph{The geometry of four-manifolds}, Oxford Mathematical Monographs,
Oxford University Press, New York, 1990.

\bibitem[DS94]{DosSal94} S.~Dostoglou and D.~A. Salamon,
\emph{Self-dual instantons and holomorphic curves}, Ann. of Math. (2)
\textbf{139} (1994), no.~3, 581--640.

\bibitem[DS07]{DosSal07} S.~Dostoglou and D.~A. Salamon,
\emph{Corrigendum to ``Self-dual instantons and holomorphic curves''},
Ann. of Math. (2) \textbf{165} (2007), no.~2, 665--673.

\bibitem[MS04]{McDuffSalamon04} D.~McDuff and D.~Salamon,
\emph{$J$-holomorphic curves and symplectic topology}, American
Mathematical Society Colloquium Publications, vol.~52, American
Mathematical Society, Providence, RI, 2004.

\bibitem[Mu\~n99a]{Munoz99a} V.~Mu\~noz, \emph{Ring structure of the
Floer cohomology of $\Sigma\times S^1$}, Topology \textbf{38} (1999),
no.~3, 517--528.

\bibitem[Mu\~n99b]{Munoz99b} V.~Mu\~noz, \emph{Quantum cohomology of the
moduli space of stable bundles over a Riemann surface}, Duke Math. J.
\textbf{98} (1999), no.~3, 525--540.

\bibitem[NS65]{NarSes65} M.~S. Narasimhan and C.~S. Seshadri,
\emph{Stable and unitary vector bundles on a compact Riemann surface},
Ann. of Math. (2) \textbf{82} (1965), no.~3, 540--567.

\bibitem[Oh93]{Oh93} Y.-G. Oh, \emph{Floer cohomology of Lagrangian
intersections and pseudo-holomorphic disks. I}, Comm. Pure Appl. Math.
\textbf{46} (1993), no.~7, 949--993.

\bibitem[Owe01]{Owens01} B.~Owens, \emph{Instantons on cylindrical
manifolds and stable bundles}, Geom. Topol. \textbf{5} (2001), 761--797.

\bibitem[Tom20]{Toma20} M.~Toma, \emph{Properness criteria for families
of coherent analytic sheaves}, Algebr. Geom. \textbf{7} (2020), no.~4,
486--502.

\bibitem[Weh05a]{Wehrheim05a} K.~Wehrheim, \emph{Anti-self-dual
instantons with Lagrangian boundary conditions. II. Bubbling}, Comm. Math.
Phys. \textbf{258} (2005), no.~2, 275--315.

\bibitem[Weh06]{Wehrheim05b} K.~Wehrheim, \emph{Energy identity for
anti-self-dual instantons on $\CC\times\Sigma$}, Math. Res. Lett.
\textbf{13} (2006), no.~1, 161--166.

\end{thebibliography}

\bigskip
\noindent{\sc Bernd Johannes Wuebben, New York, NY,}
\texttt{wuebben@gmail.com}

\end{document}
